% !TEX root = ../高等概率论.tex
% =====================================================================
%  测度延拓：半代数、代数上的测度、外测度、Carathéodory 定理、延拓定理、Lebesgue–Stieltjes 测度
% =====================================================================
\lecture{半代数与代数，代数上的测度，外测度与 Carathéodory 可测集，测度延拓定理的存在性与唯一性，Lebesgue--Stieltjes 测度与 Lebesgue 测度}
\label{lec:extension-of-measures}

上一讲解决了唯一性：测度由它在一个（适当有限的）$\pi$-类上的取值决定。这一讲解决存在性。我们以 $\R$ 上的长度为例说明问题在哪里。我们知道区间 $(a,b]$ 的长度应当是 $b-a$，我们想要的是 $\cR$ 上的一个测度 $\lambda$，使它在每个区间上取值 $b-a$。困难在于，$\cR$ 中的元素太多，我们根本写不出一般的 Borel 集长什么样，自然也无法直接写出 $\lambda$ 的定义。

一般而言，在一个“相对有限”的集合族上定义测度，要比在“相对无限”的 $\sigma$-代数上定义简单得多。有限个区间的并是可以写出来的，它们的长度也是可以直接定义的。于是我们的策略是分三步走：
\begin{enumerate}
\item 在一个结构非常简单的集合族，即左开右闭区间的全体（称为半代数）上给出集函数，验证它满足有限可加性与可数次可加性；
\item 把它延拓到有限个区间的不交并组成的代数上，这一步只是有限的组合操作；
\item 用外测度把代数上的测度延拓到生成的 $\sigma$-代数上。
\end{enumerate}
第三步是一个抽象的定理，称为 Carathéodory 延拓定理，它对任意集合上的任意代数都成立，是我们构造测度的最有效的工具。今后构造乘积测度、构造无穷维空间上的测度，走的都是这条路。

Carathéodory 定理的证明较长，其中涉及的外测度、Carathéodory 可测性等概念在课程后面几乎不再出现。教材也把这部分证明放在附录里。同学们第一次读的时候，可以先把定理的叙述和它在 Lebesgue--Stieltjes 测度上的应用弄清楚，再回过头来仔细研究证明。具体地说，第一遍可以只读下面几样东西：半代数的定义与例子；定理~\ref{thm:extension:semialgebra}、定理~\ref{thm:extension:main}、推论~\ref{cor:extension:semialgebra} 的陈述；定理~\ref{thm:extension:lebesgue-stieltjes} 的陈述、证明思路与证明；之后的全部例子与例题。标有 ★ 的证明可以跳过。

\topic{本讲用到的前文结果}

\begin{enumerate}
\item 代数与 $\sigma$-代数的定义（定义~\ref{def:set-systems:algebra}、\ref{def:set-systems:sigma-algebra}）；生成的 $\sigma$-代数的最小性（命题~\ref{prop:set-systems:generated}）。
\item $\cR$ 由 $\{(a,b]\}$ 生成（命题~\ref{prop:set-systems:borel-generators}），$\cR^d$ 由矩形 $\prod_i(a_i,b_i]$ 生成（习题~\exref{ex:set-systems:A5}）。
\item 测度的基本性质（定理~\ref{thm:measures:basic}）：单调性、次可加性、从下连续、从上连续（需要有限性）。
\item 可数可加性与连续性的关系（命题~\ref{prop:measures:additivity-continuity}）。
\item 唯一性定理（定理~\ref{thm:measures:uniqueness}）：两个测度在一个 $\pi$-类上相等，并且有一列测度有限的 $E_n\uparrow\Omega$ 属于这个 $\pi$-类，则它们在生成的 $\sigma$-代数上相等。
\item 数学分析中的 Heine--Borel 定理：闭区间 $[a,b]$（以及 $\R^d$ 中的有界闭集）的任何开覆盖都有有限子覆盖。
\item 测度的完备化（第~\ref{lec:measures} 讲作业 B），只在本讲的注记与作业 C 中用到。
\end{enumerate}

\topic{半代数与代数}

我们先说明为什么要引入半代数。左开右闭区间的全体不是代数：两个区间的并未必是区间，一个区间的补集也未必是区间。但它具有一个比代数弱的结构：两个区间的交仍是区间，一个区间的补集是有限个区间的不交并。这个结构已经足以让我们在上面谈论“有限可加性”，并且可以不费力地扩充成代数。

\begin{definition}[半代数]\label{def:extension:semialgebra}
设 $\cS\subseteq2^\Omega$ 非空。若
\begin{enumerate}
\item $S,T\in\cS$ 蕴含 $S\cap T\in\cS$；
\item 若 $S\in\cS$，则 $S^c$ 是 $\cS$ 中有限个两两不交的集合的并，
\end{enumerate}
则称 $\cS$ 为\term{半代数}（semialgebra）。
\end{definition}

由 (2)，取 $S\in\cS$，$S^c=T_1+\cdots+T_n$，$T_i\in\cS$，则 $\varnothing=S\cap T_1\in\cS$。所以半代数总含有空集。

\begin{example}[左开右闭区间]\label{ex:extension:S1}
约定 $(a,\infty]$ 表示 $(a,\infty)$。令
\[
\cS_1=\{(a,b]:\ -\infty\leqslant a<b\leqslant\infty\}\cup\{\varnothing\}.
\]
则 $\cS_1$ 是 $\R$ 上的半代数。两个区间 $(a,b]$ 与 $(a',b']$ 的交是 $(\max\{a,a'\},\min\{b,b'\}]$ 或空集，属于 $\cS_1$。$(a,b]$ 的补集是 $(-\infty,a]\cup(b,\infty]$，这是两个不交的 $\cS_1$ 中的集合的并（$a=-\infty$ 时第一项为空集，$b=\infty$ 时第二项为空集）；$\varnothing^c=(-\infty,\infty]$。
\end{example}

\begin{example}[矩形]\label{ex:extension:Sd}
令 $\cS_d$ 为空集连同一切形如 $(a_1,b_1]\times\cdots\times(a_d,b_d]$（$-\infty\leqslant a_i<b_i\leqslant\infty$）的矩形组成的集合族。$\cS_d$ 是 $\R^d$ 上的半代数。

两个矩形的交是各坐标区间之交的乘积：$\prod_i I_i\cap\prod_i J_i=\prod_i(I_i\cap J_i)$，若某个 $I_i\cap J_i=\varnothing$ 则为空集，否则每个 $I_i\cap J_i\in\cS_1$ 非空，乘积是矩形。

设 $A=I_1\times\cdots\times I_d$，$I_i\in\cS_1$ 非空。由\egref{ex:extension:S1} 的证明，$\R=I_i+J_i^{(1)}+J_i^{(2)}$，其中 $J_i^{(1)},J_i^{(2)}\in\cS_1$。记 $J_i^{(0)}=I_i$。对每个 $\varepsilon=(\varepsilon_1,\dots,\varepsilon_d)\in\{0,1,2\}^d$，令 $R_\varepsilon=J_1^{(\varepsilon_1)}\times\cdots\times J_d^{(\varepsilon_d)}$。每个 $x\in\R^d$ 的第 $i$ 个坐标恰好属于 $J_i^{(0)},J_i^{(1)},J_i^{(2)}$ 中的一个，所以 $\R^d$ 是 $3^d$ 个集合 $R_\varepsilon$ 的不交并。$A=R_{(0,\dots,0)}$，于是 $A^c$ 是其余 $3^d-1$ 个 $R_\varepsilon$ 的不交并；其中是空集的项可以去掉，非空的项都是矩形。
\end{example}

\begin{lemma}[半代数生成的代数]\label{lem:extension:algebra-from-semialgebra}
设 $\cS$ 是半代数。令 $\overline\cS$ 为 $\cS$ 中有限个两两不交的集合的并的全体。则 $\overline\cS$ 是代数，并且 $\overline\cS=\alpha(\cS)$ 是包含 $\cS$ 的最小代数。
\end{lemma}
\begin{proof}
\emph{$\overline\cS$ 对交集封闭。}设 $A=S_1+\cdots+S_m$，$B=T_1+\cdots+T_n$，$S_i,T_j\in\cS$。则
\[
A\cap B=\bigcup_{i,j}(S_i\cap T_j),
\]
其中 $S_i\cap T_j\in\cS$，并且它们两两不交：若 $(i,j)\neq(i',j')$，则 $i\neq i'$ 或 $j\neq j'$，前者时 $S_i\cap S_{i'}=\varnothing$，后者时 $T_j\cap T_{j'}=\varnothing$。所以 $A\cap B\in\overline\cS$。归纳可知 $\overline\cS$ 对有限交封闭。

\emph{$\overline\cS$ 对补集封闭。}设 $A=S_1+\cdots+S_m$。由半代数定义，每个 $S_i^c\in\overline\cS$。由 De Morgan 律，$A^c=S_1^c\cap\cdots\cap S_m^c$，由上一段属于 $\overline\cS$。

\emph{$\Omega\in\overline\cS$，并对并集封闭。}任取 $S\in\cS$，$\Omega=S\cup S^c$，而 $S^c$ 是 $\cS$ 中有限个不交集合的并，它们都与 $S$ 不交，所以 $\Omega\in\overline\cS$。$A\cup B=(A^c\cap B^c)^c\in\overline\cS$。

所以 $\overline\cS$ 是代数。它包含 $\cS$；任何包含 $\cS$ 的代数都对有限并封闭，从而包含 $\overline\cS$。
\end{proof}

例如，$\overline{\cS_1}$ 就是 $\R$ 中有限个左开右闭区间（允许无界）的不交并的全体。

\topic{代数上的测度}

\begin{definition}[代数上的测度]\label{def:extension:premeasure}
设 $\cA$ 是 $\Omega$ 上的代数。函数 $\mu:\cA\to[0,\infty]$ 若满足 $\mu(\varnothing)=0$，并且对任意两两不交的 $A_1,A_2,\dots\in\cA$，只要 $\bigcup_nA_n\in\cA$，就有 $\mu(\bigcup_nA_n)=\sum_n\mu(A_n)$，则称 $\mu$ 为代数 $\cA$ 上的\term{测度}（measure on an algebra），也称\term{预测度}（premeasure）。若存在 $A_n\in\cA$ 使 $\bigcup_nA_n=\Omega$ 且 $\mu(A_n)<\infty$，称 $\mu$ 在 $\cA$ 上是 $\sigma$-有限的。
\end{definition}

与 $\sigma$-代数上的测度相比，唯一的区别是：可数可加性只对“并集恰好还在 $\cA$ 中”的情形提出要求。取 $A_n=\varnothing$（$n>k$）可知预测度是有限可加的，于是和命题~\ref{prop:measures:additivity-continuity} 的证明一样，它是单调的。

下面的定理把半代数上的集函数延拓为代数上的测度。我们指出，从半代数到代数这一步是纯粹有限的操作，所需的只是有限可加性，唯一的技术要点是说明定义与集合的表示方法无关；真正需要分析的是可数次可加性，即下面定理中的条件 (ii)。先证明一个只用有限可加性的引理。

\begin{lemma}\label{lem:extension:finite-additive}
设 $\cS$ 是半代数，$\mu:\cS\to[0,\infty]$，$\mu(\varnothing)=0$，并且满足
\begin{enumerate}
\item[(i)] 若 $S\in\cS$ 是 $\cS$ 中有限个两两不交集合 $S_1,\dots,S_n$ 的并，则 $\mu(S)=\sum_{i=1}^n\mu(S_i)$。
\end{enumerate}
对 $A=S_1+\cdots+S_n\in\overline\cS$（$S_i\in\cS$）定义 $\bar\mu(A)=\sum_{i=1}^n\mu(S_i)$。则
\begin{enumerate}
\item $\bar\mu$ 定义合理，即与 $A$ 的表示方法无关，并且在 $\cS$ 上等于 $\mu$；
\item $\bar\mu$ 在 $\overline\cS$ 上有限可加；
\item 若 $A,B_1,\dots,B_n\in\overline\cS$ 且 $A\subseteq B_1\cup\cdots\cup B_n$，则 $\bar\mu(A)\leqslant\sum_{i=1}^n\bar\mu(B_i)$。特别地，$\bar\mu$ 单调。
\end{enumerate}
\end{lemma}
\begin{proof}
(1) 设 $A=S_1+\cdots+S_m=T_1+\cdots+T_n$ 是两种表示，$S_i,T_j\in\cS$。对每个 $i$，$S_i=S_i\cap A=\sum_j(S_i\cap T_j)$，这是 $\cS$ 中有限个不交集合的并，由 (i)，$\mu(S_i)=\sum_j\mu(S_i\cap T_j)$。同理 $\mu(T_j)=\sum_i\mu(S_i\cap T_j)$。于是
\[
\sum_i\mu(S_i)=\sum_i\sum_j\mu(S_i\cap T_j)=\sum_j\sum_i\mu(S_i\cap T_j)=\sum_j\mu(T_j).
\]
若 $A=S\in\cS$，取表示 $A=S$ 即得 $\bar\mu(S)=\mu(S)$。

(2) 设 $A=B_1+\cdots+B_n$，$A,B_k\in\overline\cS$，$B_k=\sum_jS_{kj}$，$S_{kj}\in\cS$。则 $A=\sum_{k,j}S_{kj}$ 是 $A$ 的一个表示，由 (1)，$\bar\mu(A)=\sum_{k,j}\mu(S_{kj})=\sum_k\bar\mu(B_k)$。

(3) 先看 $n=1$：若 $A\subseteq B$，则 $B=A+(B\cap A^c)$，而 $B\cap A^c\in\overline\cS$（$\overline\cS$ 是代数），由 (2)，$\bar\mu(B)=\bar\mu(A)+\bar\mu(B\cap A^c)\geqslant\bar\mu(A)$。一般情形，令 $F_1=B_1$，$F_k=B_k\cap B_1^c\cap\cdots\cap B_{k-1}^c$，则 $F_k\in\overline\cS$ 两两不交，$F_k\subseteq B_k$，$\bigcup_kF_k=\bigcup_kB_k$。于是 $A=\sum_k(A\cap F_k)$，由 (2) 与单调性，
\[
\bar\mu(A)=\sum_k\bar\mu(A\cap F_k)\leqslant\sum_k\bar\mu(F_k)\leqslant\sum_k\bar\mu(B_k).\qedhere
\]
\end{proof}

\begin{example}[半代数上的集函数]
在半代数 $\cS_1$ 上令 $\mu((a,b])=b-a$。条件 (i) 说的是：若 $(a,b]$ 被切成有限个不交的小区间，长度相加等于总长度。比如 $(0,3]=(0,1]+(1,2]+(2,3]$，$3=1+1+1$。条件 (ii) 说的是：若 $(a,b]$ 被切成\emph{可数}个不交的小区间，长度之和不小于总长度。比如 $(0,1]=\sum_{n\geqslant1}(2^{-n},2^{-n+1}]$，$\sum_n2^{-n}=1$。(ii) 看起来理所当然，但它的证明需要实数的完备性（紧性）：在有理数集 $\Q$ 上，$\Q\cap(0,1]$ 可以被可数个总长度任意小的有理区间盖住（例题，第~\pageref{we:measures:rationals} 页），“长度”就没有可数次可加性。所以条件 (ii) 是构造中唯一需要动脑筋的地方。
\end{example}

\begin{theorem}[从半代数到代数]\label{thm:extension:semialgebra}
设 $\cS$ 是半代数，$\mu:\cS\to[0,\infty]$，$\mu(\varnothing)=0$，满足引理~\ref{lem:extension:finite-additive} 中的条件 (i)，以及
\begin{enumerate}
\item[(ii)] 若 $S\in\cS$ 是 $\cS$ 中可数个两两不交集合 $S_1,S_2,\dots$ 的并，则 $\mu(S)\leqslant\sum_{i=1}^\infty\mu(S_i)$。
\end{enumerate}
则引理~\ref{lem:extension:finite-additive} 中的 $\bar\mu$ 是代数 $\overline\cS$ 上的测度，并且是 $\mu$ 在 $\overline\cS$ 上唯一的有限可加延拓。
\end{theorem}
\begin{proof}
唯一性：任何有限可加的延拓在 $A=S_1+\cdots+S_n$ 上都必须取值 $\sum_i\mu(S_i)$。

可数可加性：设 $A,B_1,B_2,\dots\in\overline\cS$，$A=\sum_{i\geqslant1}B_i$。

\emph{“$\geqslant$”。}对每个 $n$，$A\supseteq B_1+\cdots+B_n$，由引理~\ref{lem:extension:finite-additive} 的单调性与有限可加性，$\bar\mu(A)\geqslant\sum_{i\leqslant n}\bar\mu(B_i)$。令 $n\to\infty$ 得 $\bar\mu(A)\geqslant\sum_i\bar\mu(B_i)$。

\emph{“$\leqslant$”。}写 $B_i=\sum_jS_{ij}$（有限和），$A=\sum_{k=1}^mT_k$，$S_{ij},T_k\in\cS$。对每个 $k$，$T_k=T_k\cap A=\sum_{i,j}(T_k\cap S_{ij})$ 是 $\cS$ 中可数个不交集合的并，由 (ii)，
\[
\mu(T_k)\leqslant\sum_{i,j}\mu(T_k\cap S_{ij}).
\]
另一方面，由于 $S_{ij}\subseteq A$，$S_{ij}=\sum_k(T_k\cap S_{ij})$ 是有限不交并，由 (i)，$\mu(S_{ij})=\sum_k\mu(T_k\cap S_{ij})$。对 $k$ 求和并交换非负级数的求和次序，
\[
\bar\mu(A)=\sum_k\mu(T_k)\leqslant\sum_k\sum_{i,j}\mu(T_k\cap S_{ij})=\sum_{i,j}\sum_k\mu(T_k\cap S_{ij})=\sum_{i,j}\mu(S_{ij})=\sum_i\bar\mu(B_i).\qedhere
\]
\end{proof}

\topic{外测度与 Carathéodory 定理}

现在我们有了代数 $\cA$ 上的测度 $\mu$，要把它延拓到 $\sigma(\cA)$ 上去。想法其实很朴素，同学们在数学分析中学 Riemann 积分的时候已经见过：对任意集合 $E$，用 $\cA$ 中的集合从外面把它盖住，盖住所用的集合的测度之和总是不小于 $E$ 的“大小”的，对一切覆盖取下确界，就得到 $E$ 的一个合理的“大小”，称为外测度。外测度对一切子集都有定义，但它一般不是可加的。Carathéodory 找到了一个条件，把使外测度可加的那些集合挑出来，它们恰好组成一个 $\sigma$-代数，并且包含 $\cA$。

\begin{definition}[外测度]\label{def:extension:outer}
函数 $\mu^*:2^\Omega\to[0,\infty]$ 若满足
\begin{enumerate}
\item $\mu^*(\varnothing)=0$；
\item 单调性：$E\subseteq F$ 蕴含 $\mu^*(E)\leqslant\mu^*(F)$；
\item 可数次可加性：$E\subseteq\bigcup_nE_n$ 蕴含 $\mu^*(E)\leqslant\sum_n\mu^*(E_n)$，
\end{enumerate}
则称 $\mu^*$ 为 $\Omega$ 上的\term{外测度}（outer measure）。
\end{definition}

\begin{proposition}[由代数上的测度诱导的外测度]\label{prop:extension:induced-outer}
设 $\mu$ 是代数 $\cA$ 上的测度。对 $E\subseteq\Omega$ 令
\[
\mu^*(E)=\inf\Bigl\{\sum_{n=1}^\infty\mu(A_n):\ A_n\in\cA,\ E\subseteq\bigcup_{n=1}^\infty A_n\Bigr\}.
\]
则 $\mu^*$ 是外测度，并且对 $A\in\cA$ 有 $\mu^*(A)=\mu(A)$。
\end{proposition}
\begin{proof}
由于 $\Omega\in\cA$，每个 $E$ 至少有覆盖 $A_1=\Omega$，$A_n=\varnothing$（$n\geqslant2$），所以下确界是对非空集合取的。

(1) 取 $A_n=\varnothing$ 得 $\mu^*(\varnothing)=0$。

(2) 若 $E\subseteq F$，则 $F$ 的每个覆盖都是 $E$ 的覆盖，$E$ 的下确界是对更大的集合取的，所以 $\mu^*(E)\leqslant\mu^*(F)$。

(3) 设 $E\subseteq\bigcup_nE_n$。若某个 $\mu^*(E_n)=\infty$，结论显然。否则任取 $\varepsilon>0$，对每个 $n$，由下确界的定义取 $A_{n,k}\in\cA$ 使 $E_n\subseteq\bigcup_kA_{n,k}$ 且 $\sum_k\mu(A_{n,k})\leqslant\mu^*(E_n)+\varepsilon2^{-n}$。则 $\{A_{n,k}\}_{n,k}$ 是 $E$ 的可数覆盖，于是
\[
\mu^*(E)\leqslant\sum_n\sum_k\mu(A_{n,k})\leqslant\sum_n\bigl(\mu^*(E_n)+\varepsilon2^{-n}\bigr)=\sum_n\mu^*(E_n)+\varepsilon.
\]
由 $\varepsilon$ 的任意性得结论。

最后设 $A\in\cA$。取覆盖 $A_1=A$，其余为空集，得 $\mu^*(A)\leqslant\mu(A)$。反之，设 $A\subseteq\bigcup_nA_n$，$A_n\in\cA$。令 $B_1=A\cap A_1$，$B_n=A\cap A_n\cap A_1^c\cap\cdots\cap A_{n-1}^c$，则 $B_n\in\cA$ 两两不交，$B_n\subseteq A_n$，$\bigcup_nB_n=A\in\cA$。由 $\mu$ 在代数上的可数可加性与单调性，$\mu(A)=\sum_n\mu(B_n)\leqslant\sum_n\mu(A_n)$。对一切覆盖取下确界得 $\mu(A)\leqslant\mu^*(A)$。
\end{proof}

\begin{example}[覆盖必须是可数的]
我们强调，外测度定义中的覆盖 $\{A_n\}$ 是可数个集合，不能换成有限个。考虑 $[0,1]$ 中有限个区间的并组成的代数，其上的测度是区间的长度。如果只允许用有限个区间来盖住 $\Q\cap[0,1]$，那么这有限个区间的并是闭集（把区间都换成闭区间不改变长度），它包含 $\Q\cap[0,1]$，从而包含它的闭包 $[0,1]$，长度之和至少为 $1$。可是如果允许可数个区间，把第 $n$ 个有理数 $q_n$ 用长度为 $\varepsilon2^{-n}$ 的区间盖住，长度之和不超过 $\varepsilon$，所以 $\Q\cap[0,1]$ 的外测度是 $0$。前一种“有限覆盖”的做法对应的是 Jordan 测度，它与 Riemann 积分相配；Lebesgue 的想法正是把有限换成可数。$\Q\cap[0,1]$ 的示性函数不是 Riemann 可积的，而在 Lebesgue 的理论中它的积分就是 $0$。
\end{example}

外测度还有一个直观的意义。若 $\nu$ 是 $\sigma(\cA)$ 上的测度，并且在 $\cA$ 上与 $\mu$ 相等，那么对 $E\in\sigma(\cA)$ 和 $E$ 的任何覆盖 $\{A_n\}$，由次可加性，$\nu(E)\leqslant\sum_n\nu(A_n)=\sum_n\mu(A_n)$，从而 $\nu(E)\leqslant\mu^*(E)$。所以外测度是一切可能的延拓在 $E$ 上所能取到的值的上界。我们希望找到一些集合，对它们而言这个上界是“准确的”。Carathéodory 把这个想法表述成下面的定义。

\begin{definition}[Carathéodory 可测]\label{def:extension:caratheodory}
设 $\mu^*$ 是 $\Omega$ 上的外测度。若 $E\subseteq\Omega$ 满足
\[
\mu^*(F)=\mu^*(F\cap E)+\mu^*(F\cap E^c)\qquad\text{对一切}\ F\subseteq\Omega,
\]
则称 $E$ 是 \term{$\mu^*$-可测的}（$\mu^*$-measurable）。$\mu^*$-可测集的全体记作 $\cM(\mu^*)$。
\end{definition}

定义中的 $F$ 称为试验集。$E$ 是 $\mu^*$-可测的，意思是用 $E$ 把任何一个集合 $F$ 切成两块，两块的外测度之和恰好等于 $F$ 的外测度。由次可加性，$\mu^*(F)\leqslant\mu^*(F\cap E)+\mu^*(F\cap E^c)$ 总成立，所以要验证 $E$ 可测，只需验证
\begin{equation}\label{eq:extension:caratheodory-half}
\mu^*(F)\geqslant\mu^*(F\cap E)+\mu^*(F\cap E^c)\qquad\text{对一切}\ F\ \text{满足}\ \mu^*(F)<\infty
\end{equation}
（$\mu^*(F)=\infty$ 时不等式自动成立）。

\begin{remark}[证明思路]
要证明 $\cM$ 是 $\sigma$-代数、$\mu^*$ 在 $\cM$ 上可数可加。先证 $\cM$ 对补集和有限并封闭（这只用到定义）；再证有限个不交可测集上 $\mu^*$ 可加，并且“可加”对任意试验集 $G$ 截下来的部分也成立；最后取极限，同时得到可数并的可测性与可数可加性。全部证明只用到外测度的三条性质（定义~\ref{def:extension:outer}）。
\end{remark}

\begin{theorem}[Carathéodory 定理]\label{thm:extension:caratheodory}
设 $\mu^*$ 是 $\Omega$ 上的外测度。则 $\cM=\cM(\mu^*)$ 是 $\sigma$-代数，$\mu^*$ 在 $\cM$ 上的限制是测度。此外，若 $\mu^*(N)=0$，则 $N\in\cM$。
\end{theorem}
\begin{proof}[证明★]
分五步。

\emph{(a) $\Omega\in\cM$，$\cM$ 对补集封闭。}定义中 $E$ 与 $E^c$ 的地位对称，所以 $E\in\cM$ 蕴含 $E^c\in\cM$。$\mu^*(F\cap\Omega)+\mu^*(F\cap\varnothing)=\mu^*(F)$，故 $\Omega\in\cM$。

\emph{(b) $\cM$ 对有限并封闭。}设 $E_1,E_2\in\cM$，$G\subseteq\Omega$。用 $E_1$ 可测（试验集 $G$）、$E_2$ 可测（试验集 $G\cap E_1^c$），
\[
\mu^*(G)=\mu^*(G\cap E_1)+\mu^*(G\cap E_1^c)=\mu^*(G\cap E_1)+\mu^*(G\cap E_1^c\cap E_2)+\mu^*(G\cap E_1^c\cap E_2^c).
\]
由于 $G\cap(E_1\cup E_2)=(G\cap E_1)\cup(G\cap E_1^c\cap E_2)$，由次可加性，前两项之和不小于 $\mu^*(G\cap(E_1\cup E_2))$；而第三项是 $\mu^*(G\cap(E_1\cup E_2)^c)$。所以
\[
\mu^*(G)\geqslant\mu^*(G\cap(E_1\cup E_2))+\mu^*(G\cap(E_1\cup E_2)^c),
\]
即 $E_1\cup E_2\in\cM$。归纳得有限并封闭。结合 (a)，$\cM$ 是代数，从而对有限交与差集也封闭。

\emph{(c) 若 $E_1,\dots,E_n\in\cM$ 两两不交，则对任意 $G$，$\mu^*(G\cap\bigcup_{i\leqslant n}E_i)=\sum_{i\leqslant n}\mu^*(G\cap E_i)$。}对 $n$ 归纳，$n=1$ 显然。令 $H_n=E_1\cup\cdots\cup E_n$。用 $E_n$ 可测，试验集取 $G\cap H_n$：由于 $H_n\cap E_n=E_n$，$H_n\cap E_n^c=H_{n-1}$（不交性），
\[
\mu^*(G\cap H_n)=\mu^*(G\cap H_n\cap E_n)+\mu^*(G\cap H_n\cap E_n^c)=\mu^*(G\cap E_n)+\mu^*(G\cap H_{n-1}),
\]
由归纳假设即得。

\emph{(d) $\cM$ 对可数并封闭。}设 $E_i\in\cM$，$E=\bigcup_iE_i$。令 $E_1'=E_1$，$E_i'=E_i\setminus(E_1\cup\cdots\cup E_{i-1})$，由 (b) 知 $E_i'\in\cM$，它们两两不交且并为 $E$。所以不妨设 $E_i$ 两两不交。令 $H_n=E_1\cup\cdots\cup E_n\in\cM$。对任意 $G$，由 $H_n$ 可测、$H_n^c\supseteq E^c$ 与单调性、以及 (c)，
\[
\mu^*(G)=\mu^*(G\cap H_n)+\mu^*(G\cap H_n^c)\geqslant\sum_{i=1}^n\mu^*(G\cap E_i)+\mu^*(G\cap E^c).
\]
令 $n\to\infty$，再用次可加性（$G\cap E=\bigcup_i(G\cap E_i)$），
\begin{equation}\label{eq:extension:countable}
\mu^*(G)\geqslant\sum_{i=1}^\infty\mu^*(G\cap E_i)+\mu^*(G\cap E^c)\geqslant\mu^*(G\cap E)+\mu^*(G\cap E^c).
\end{equation}
所以 $E\in\cM$。

\emph{(e) $\mu^*$ 在 $\cM$ 上可数可加。}设 $E_i\in\cM$ 两两不交，$E=\bigcup_iE_i$。在 \eqref{eq:extension:countable} 中取 $G=E$，得 $\mu^*(E)\geqslant\sum_i\mu^*(E_i)$；反向不等式是次可加性。又 $\mu^*(\varnothing)=0$，故 $\mu^*|_\cM$ 是测度。

最后，设 $\mu^*(N)=0$。对任意 $F$，由单调性 $\mu^*(F\cap N)\leqslant\mu^*(N)=0$，$\mu^*(F\cap N^c)\leqslant\mu^*(F)$，所以 \eqref{eq:extension:caratheodory-half} 成立，$N\in\cM$。
\end{proof}

\begin{remark}[Carathéodory 条件的直观]
Carathéodory 的条件初看起来很不自然，我们来解释它为什么是合理的。以 $\R$ 上的长度为例，设 $E\subseteq I=(a,b]$。在条件中取试验集 $F=I$，得到
\[
\lambda^*(E)=(b-a)-\lambda^*(I\setminus E).
\]
左边是用区间从外面盖住 $E$ 得到的大小，右边是用 $I$ 的长度减去 $E$ 在 $I$ 中的补集的外测度，这可以理解为从里面逼近 $E$ 得到的大小（称为内测度）。所以对这样的 $E$，Carathéodory 条件说的就是“外测度等于内测度”，这正是 Lebesgue 当年定义可测集的方式。Carathéodory 的高明之处在于，他不需要定义内测度，而是要求等式对一切试验集 $F$ 都成立，这样一来，条件就只涉及外测度本身，对任意集合上的任意外测度都有意义，证明也变得非常整齐。
\end{remark}

\begin{lemma}\label{lem:extension:algebra-measurable}
设 $\mu$ 是代数 $\cA$ 上的测度，$\mu^*$ 是它诱导的外测度。则 $\cA\subseteq\cM(\mu^*)$。
\end{lemma}
\begin{proof}[证明★]
设 $A\in\cA$，要对 $\mu^*(F)<\infty$ 的 $F$ 验证 \eqref{eq:extension:caratheodory-half}。任取 $\varepsilon>0$，取 $B_n\in\cA$ 使 $F\subseteq\bigcup_nB_n$ 且 $\sum_n\mu(B_n)\leqslant\mu^*(F)+\varepsilon$。对每个 $n$，$B_n=(B_n\cap A)+(B_n\cap A^c)$ 是 $\cA$ 中两个不交集合的并，由 $\mu$ 的有限可加性以及 $\mu=\mu^*$ 在 $\cA$ 上成立（命题~\ref{prop:extension:induced-outer}），
\[
\mu(B_n)=\mu^*(B_n\cap A)+\mu^*(B_n\cap A^c).
\]
对 $n$ 求和，并由 $F\cap A\subseteq\bigcup_n(B_n\cap A)$，$F\cap A^c\subseteq\bigcup_n(B_n\cap A^c)$ 与次可加性，
\[
\mu^*(F)+\varepsilon\geqslant\sum_n\mu^*(B_n\cap A)+\sum_n\mu^*(B_n\cap A^c)\geqslant\mu^*(F\cap A)+\mu^*(F\cap A^c).
\]
由 $\varepsilon$ 的任意性得 \eqref{eq:extension:caratheodory-half}。
\end{proof}

\topic{测度延拓定理}

\begin{theorem}[Carathéodory 测度延拓定理]\label{thm:extension:main}
设 $\mu$ 是代数 $\cA$ 上的测度。
\begin{enumerate}
\item \emph{存在性}：存在 $\sigma(\cA)$ 上的测度 $\nu$，使 $\nu(A)=\mu(A)$ 对一切 $A\in\cA$ 成立。可以取 $\nu=\mu^*|_{\sigma(\cA)}$，其中 $\mu^*$ 是 $\mu$ 诱导的外测度。
\item \emph{唯一性}：若 $\mu$ 在 $\cA$ 上 $\sigma$-有限，则这样的 $\nu$ 是唯一的。
\end{enumerate}
\end{theorem}
\begin{proof}
(1) 由引理~\ref{lem:extension:algebra-measurable}，$\cA\subseteq\cM(\mu^*)$；由定理~\ref{thm:extension:caratheodory}，$\cM(\mu^*)$ 是 $\sigma$-代数，故 $\sigma(\cA)\subseteq\cM(\mu^*)$。$\mu^*$ 在 $\cM(\mu^*)$ 上是测度，限制到子 $\sigma$-代数 $\sigma(\cA)$ 上仍是测度（可数可加性对更少的集合列成立）。由命题~\ref{prop:extension:induced-outer}，它在 $\cA$ 上等于 $\mu$。

(2) 设 $\nu_1,\nu_2$ 是两个延拓。$\cA$ 是代数，从而是 $\pi$-类。由 $\sigma$-有限性，存在 $A_n\in\cA$，$\bigcup_nA_n=\Omega$，$\mu(A_n)<\infty$；令 $E_n=A_1\cup\cdots\cup A_n\in\cA$，则 $E_n\uparrow\Omega$ 且 $\mu(E_n)\leqslant\sum_{k\leqslant n}\mu(A_k)<\infty$。由唯一性定理（定理~\ref{thm:measures:uniqueness}），$\nu_1=\nu_2$ 在 $\sigma(\cA)$ 上成立。
\end{proof}

把定理~\ref{thm:extension:semialgebra} 与定理~\ref{thm:extension:main} 合起来，就得到构造测度时实际使用的形式。

\begin{corollary}[从半代数延拓]\label{cor:extension:semialgebra}
设 $\cS$ 是半代数，$\mu:\cS\to[0,\infty]$ 满足 $\mu(\varnothing)=0$ 以及定理~\ref{thm:extension:semialgebra} 的条件 (i)(ii)。则存在 $\sigma(\cS)$ 上的测度 $\nu$，使 $\nu=\mu$ 在 $\cS$ 上成立。若存在 $S_n\in\cS$ 使 $\bigcup_nS_n=\Omega$ 且 $\mu(S_n)<\infty$，则 $\nu$ 唯一。
\end{corollary}
\begin{proof}
由定理~\ref{thm:extension:semialgebra}，$\mu$ 延拓为代数 $\overline\cS$ 上的测度 $\bar\mu$。由引理~\ref{lem:extension:algebra-from-semialgebra}，$\sigma(\overline\cS)=\sigma(\cS)$（$\cS\subseteq\overline\cS\subseteq\sigma(\cS)$）。由定理~\ref{thm:extension:main}，$\bar\mu$ 延拓为 $\sigma(\cS)$ 上的测度。唯一性：任何在 $\cS$ 上等于 $\mu$ 的测度在 $\overline\cS$ 上都等于 $\bar\mu$（有限可加性），而 $\bar\mu$ 在 $\overline\cS$ 上 $\sigma$-有限，由定理~\ref{thm:extension:main}(2) 即得。
\end{proof}

\begin{remark}
我们把这一路走下来的过程总结一下。要在 $\sigma(\cS)$ 上构造一个测度，只需要在半代数 $\cS$ 上给出集函数，并验证两件事：有限可加性 (i) 与可数次可加性 (ii)。其余的一切，即延拓到代数、定义外测度、挑出可测集、验证可数可加性，都由一般的理论自动完成，与具体的问题无关。所以在具体问题中，我们的全部工作都集中在验证 (i) 和 (ii) 上，而 (ii) 往往是唯一需要动脑筋的地方。下一节的 Lebesgue--Stieltjes 测度是第一个例子。
\end{remark}

\begin{example}[没有 $\sigma$-有限性时延拓可以不唯一]
在 $\R$ 上取代数 $\overline{\cS_1}$，令 $\mu(\varnothing)=0$，$\mu(A)=\infty$（$A\neq\varnothing$）。$\mu$ 是 $\overline{\cS_1}$ 上的测度：一列不交集合的并非空，当且仅当其中某一项非空。以下三个 $\cR$ 上的测度都是它的延拓：$\nu_1(B)=\#(B\cap\Q)$，$\nu_2(B)=2\#(B\cap\Q)$，以及 $\nu_3(B)=\infty\cdot\one_{\{B\neq\varnothing\}}$（$\nu_3$ 是测度：可数个不交集合的并非空当且仅当某一项非空）。事实上，非空的 $\overline{\cS_1}$ 中的集合含有非退化区间，从而含有无穷多个有理数。
\end{example}

\begin{remark}
定理~\ref{thm:extension:main} 的证明给出的其实是一个比 $\sigma(\cA)$ 更大的 $\sigma$-代数 $\cM(\mu^*)$ 上的测度。当 $\mu$ 是 $\sigma$-有限时，$(\Omega,\cM(\mu^*),\mu^*)$ 恰好是 $(\Omega,\sigma(\cA),\mu^*)$ 的完备化（上一讲作业 B），证明见本讲作业 C。对 $\R$ 上的长度，$\cM(\lambda^*)$ 称为 \term{Lebesgue 可测集}（Lebesgue measurable sets）的全体，它严格大于 $\cR$（思考题 E），但仍不是 $2^\R$（作业 D）。在概率论中，我们几乎总是在 $\sigma(\cA)$ 上工作。
\end{remark}

\topic{Lebesgue--Stieltjes 测度与 Lebesgue 测度}

现在把延拓定理用到 $\R$ 上。我们不只构造长度，而是一次性地构造一大类测度：给定一个单调不减的函数 $F$，我们想让区间 $(a,b]$ 的“质量”等于 $F(b)-F(a)$。$F(x)=x$ 时这就是长度；以后会看到，$F$ 取为分布函数时，这就是随机变量的分布。验证有限可加性只是区间端点的排列组合，验证可数次可加性则要用到闭区间的紧性和 $F$ 的右连续性，这是整个构造中唯一真正用到 $\R$ 的拓扑性质的地方。

\begin{definition}[Stieltjes 测度函数]\label{def:extension:stieltjes}
函数 $F:\R\to\R$ 若满足
\begin{enumerate}
\item $F$ 单调不减；
\item $F$ 右连续：对一切 $x$，$\lim_{y\downarrow x}F(y)=F(x)$，
\end{enumerate}
则称 $F$ 为 \term{Stieltjes 测度函数}（Stieltjes measure function）。
\end{definition}

由单调性，极限 $F(\infty)=\lim_{x\to\infty}F(x)\in(-\infty,\infty]$ 与 $F(-\infty)=\lim_{x\to-\infty}F(x)\in[-\infty,\infty)$ 存在。对 $-\infty\leqslant a<b\leqslant\infty$，$F(b)-F(a)$ 总有意义（不会出现 $\infty-\infty$），取值于 $[0,\infty]$。

\begin{remark}[证明思路]
由推论~\ref{cor:extension:semialgebra}，只需在半代数 $\cS_1$ 上验证条件 (i) 与 (ii)。(i) 是区间端点的排列：一个区间被切成有限个小区间，小区间首尾相接，$F$ 的增量逐项相消。(ii) 分三步：先把覆盖区间都截到 $(a,b]$ 里；再把 $(a,b]$ 缩成闭区间 $[a+\delta,b]$、把每个覆盖区间放大成开区间，用紧性化成有限覆盖，由 (i) 得到的有限次可加性估计；最后令放大缩小的量趋于零，这里用到 $F$ 的右连续性。唯一性直接由唯一性定理得到。
\end{remark}

\begin{theorem}[Lebesgue--Stieltjes 测度]\label{thm:extension:lebesgue-stieltjes}
设 $F$ 是 Stieltjes 测度函数。则 $(\R,\cR)$ 上存在唯一的测度 $\mu_F$，使
\[
\mu_F((a,b])=F(b)-F(a)\qquad\text{对一切}\ -\infty<a<b<\infty
\]
成立。$\mu_F$ 称为 $F$ 对应的 \term{Lebesgue--Stieltjes 测度}（Lebesgue--Stieltjes measure）。$F(x)=x$ 对应的测度记作 $\lambda$，称为 \term{Lebesgue 测度}（Lebesgue measure）。
\end{theorem}
\begin{proof}
在半代数 $\cS_1$ 上令 $\mu(\varnothing)=0$，$\mu((a,b])=F(b)-F(a)$（$-\infty\leqslant a<b\leqslant\infty$）。我们验证推论~\ref{cor:extension:semialgebra} 的条件。

\emph{条件 (i)：有限可加性。}设 $(a,b]=\sum_{i=1}^n(a_i,b_i]$，各项非空（空集的项贡献 $0$，可以去掉）。由于这些区间两两不交，它们的左端点互不相同（若 $a_i=a_j$，则两个区间都含有 $a_i$ 右侧附近的点）。按左端点从小到大重新编号，使 $a_1<a_2<\cdots<a_n$。
\begin{itemize}
\item $b_i\leqslant a_{i+1}$：否则 $a_{i+1}<b_i$，而 $a_i<a_{i+1}$，于是 $(a_{i+1},\min\{b_i,b_{i+1}\}]$ 非空，且同时含于第 $i$ 个和第 $i+1$ 个区间，与不交矛盾。于是 $b_1\leqslant a_2<b_2\leqslant a_3<\cdots$，特别地 $b_n=\max_ib_i$。
\item $a_1=a$，$b_n=b$：并集中的点都大于 $a_1$，并且 $a_1$ 右侧附近的点属于并集，所以并集的下确界是 $a_1$；$(a,b]$ 的下确界是 $a$，故 $a_1=a$。同理比较上确界（并集的上确界是 $b_n$）得 $b_n=b$。
\item $b_i=a_{i+1}$：若 $b_i<a_{i+1}$，则 $(b_i,a_{i+1}]\subseteq(a_1,b_n]=(a,b]$，但其中的点大于 $b_j$（$j\leqslant i$）且不超过 $a_j$（$j\geqslant i+1$），不属于任何一个区间，矛盾。
\end{itemize}
于是 $\sum_i(F(b_i)-F(a_i))=F(b_n)-F(a_1)=F(b)-F(a)$（逐项相消）。

条件 (i) 成立后，引理~\ref{lem:extension:finite-additive} 的结论可用，特别是有限次可加性：若 $(c,d]\subseteq\bigcup_{j=1}^J(c_j,d_j]$，则 $F(d)-F(c)\leqslant\sum_j(F(d_j)-F(c_j))$。

\emph{条件 (ii)：可数次可加性。}我们证明更强的结论：若 $(a,b]\subseteq\bigcup_{i\geqslant1}(a_i,b_i]$（不要求不交），则
\begin{equation}\label{eq:extension:countable-subadd}
F(b)-F(a)\leqslant\sum_{i\geqslant1}\bigl(F(b_i)-F(a_i)\bigr).
\end{equation}

\emph{情形一：$a,b$ 有限。}把 $(a_i,b_i]$ 换成 $(a_i,b_i]\cap(a,b]=(a_i',b_i']$（为空集时去掉），覆盖关系不变，并且由 $F$ 单调，$F(b_i')-F(a_i')\leqslant F(b_i)-F(a_i)$。所以不妨设 $a\leqslant a_i<b_i\leqslant b$，全部有限。任取 $\varepsilon>0$。由右连续性，取 $\delta\in(0,b-a)$ 使 $F(a+\delta)<F(a)+\varepsilon$，对每个 $i$ 取 $\eta_i>0$ 使 $F(b_i+\eta_i)<F(b_i)+\varepsilon2^{-i}$。开区间族 $\{(a_i,b_i+\eta_i)\}_i$ 覆盖闭区间 $[a+\delta,b]$（因为 $[a+\delta,b]\subseteq(a,b]\subseteq\bigcup_i(a_i,b_i]$）。由 Heine--Borel 定理，存在有限个 $i_1,\dots,i_J$ 使 $[a+\delta,b]\subseteq\bigcup_{j\leqslant J}(a_{i_j},b_{i_j}+\eta_{i_j})$，从而 $(a+\delta,b]\subseteq\bigcup_{j\leqslant J}(a_{i_j},b_{i_j}+\eta_{i_j}]$。由有限次可加性，
\[
F(b)-F(a+\delta)\leqslant\sum_{j=1}^J\bigl(F(b_{i_j}+\eta_{i_j})-F(a_{i_j})\bigr)\leqslant\sum_{i=1}^\infty\bigl(F(b_i)-F(a_i)+\varepsilon2^{-i}\bigr).
\]
于是 $F(b)-F(a)<F(b)-F(a+\delta)+\varepsilon\leqslant\sum_i(F(b_i)-F(a_i))+2\varepsilon$。由 $\varepsilon$ 的任意性得 \eqref{eq:extension:countable-subadd}。

\emph{情形二：$a=-\infty$ 或 $b=\infty$。}任取有限的 $A<B$ 使 $(A,B]\subseteq(a,b]$，由情形一，$F(B)-F(A)\leqslant\sum_i(F(b_i)-F(a_i))$。令 $A\downarrow a$，$B\uparrow b$（$a$ 有限时取 $A=a$，$b$ 有限时取 $B=b$），由 $F(\pm\infty)$ 的定义，左边趋于 $F(b)-F(a)$。

\emph{$\sigma$-有限与结论。}$\R=\bigcup_n(-n,n]$，$\mu((-n,n])=F(n)-F(-n)<\infty$。由推论~\ref{cor:extension:semialgebra}，存在 $\sigma(\cS_1)=\cR$（命题~\ref{prop:set-systems:borel-generators}）上的测度 $\mu_F$，它在 $\cS_1$ 上等于 $\mu$，从而满足定理的等式。

\emph{唯一性。}设 $\nu$ 是 $\cR$ 上满足 $\nu((a,b])=F(b)-F(a)$（$a,b$ 有限）的测度。$\cP=\{(a,b]:-\infty<a<b<\infty\}\cup\{\varnothing\}$ 是生成 $\cR$ 的 $\pi$-类，$(-n,n]\in\cP$，$(-n,n]\uparrow\R$，$\nu((-n,n])<\infty$。由唯一性定理（定理~\ref{thm:measures:uniqueness}），$\nu=\mu_F$。
\end{proof}

\begin{remark}
我们回顾一下这个证明中紧性是怎样起作用的。可数次可加性要求：若 $(a,b]$ 被可数个区间盖住，则长度之和不小于 $b-a$。对有限个区间，这是显然的；困难在于可数个。紧性把可数覆盖化成有限覆盖，但紧性只对闭区间和开覆盖成立，所以我们付出了两个 $\varepsilon$ 的代价：把 $(a,b]$ 缩成闭区间 $[a+\delta,b]$，把每个 $(a_i,b_i]$ 放大成开区间 $(a_i,b_i+\eta_i)$。右连续性保证了这两处代价可以任意小。这个“缩小一点、放大一点、用紧性、再令 $\varepsilon\to0$”的论证，同学们应该熟记，它在测度论中反复出现。
\end{remark}

\begin{proposition}[Lebesgue 测度的基本性质]\label{prop:extension:lebesgue}
\begin{enumerate}
\item 对任意 $x\in\R$，$\lambda(\{x\})=0$；从而对 $a<b$，$\lambda([a,b])=\lambda((a,b))=\lambda([a,b))=b-a$，并且每个可数集都是 $\lambda$-零测集。
\item \term{平移不变性}（translation invariance）：对 $A\in\cR$ 与 $x\in\R$，$A+x=\{a+x:a\in A\}\in\cR$，且 $\lambda(A+x)=\lambda(A)$。
\end{enumerate}
\end{proposition}
\begin{proof}
(1) $(x-1/n,x]\downarrow\{x\}$，$\lambda((x-1,x])=1<\infty$，由从上连续，$\lambda(\{x\})=\lim_n1/n=0$。于是 $\lambda([a,b])=\lambda(\{a\})+\lambda((a,b])=b-a$，其余同理。可数集是可数个单点集的并，由可数可加性测度为 $0$。

(2) 令 $\cG=\{A\in\cR: A+x\in\cR\}$。$\cG$ 是 $\sigma$-代数：$\R+x=\R$；$(A+x)^c=A^c+x$；$(\bigcup_nA_n)+x=\bigcup_n(A_n+x)$。$\cG$ 包含一切 $(a,b]$，因为 $(a,b]+x=(a+x,b+x]$。所以 $\cG\supseteq\sigma(\cS_1)=\cR$。现在令 $\nu(A)=\lambda(A+x)$，由刚才的两个集合等式，$\nu$ 是 $\cR$ 上的测度，且 $\nu((a,b])=(b+x)-(a+x)=b-a$。由定理~\ref{thm:extension:lebesgue-stieltjes} 的唯一性，$\nu=\lambda$。
\end{proof}

\begin{remark}
平移不变性在本质上刻画了 Lebesgue 测度：若 $\mu$ 是 $\cR$ 上的测度，在平移下不变，并且 $\mu((0,1])=1$，那么 $\mu=\lambda$。证明的想法很几何：$(0,1]$ 是 $n$ 个平移后互相重合的小区间 $(\frac{k-1}n,\frac kn]$ 的不交并，所以每个小区间的测度是 $1/n$；于是端点为有理数的区间的测度等于长度；再用单调性和连续性推到一切区间，最后用唯一性（习题~\exref{ex:extension:A10}）。这说明长度这个概念并不是人为规定的：只要求“平移不变”和“单位区间长度为 $1$”，长度就被唯一确定了。
\end{remark}

\begin{example}[几个 Lebesgue--Stieltjes 测度]\label{ex:extension:LS-examples}
\begin{enumerate}
\item $F=\one_{[0,\infty)}$：单调不减，右连续。对 $a<b$，$F(b)-F(a)=1$ 当且仅当 $a<0\leqslant b$，即 $0\in(a,b]$。所以 $\mu_F((a,b])=\delta_0((a,b])$，由唯一性，$\mu_F=\delta_0$。注意若取左连续的 $\one_{(0,\infty)}$，它不是 Stieltjes 测度函数，这就是为什么我们约定使用左开右闭区间与右连续函数：$(a,b_n]\downarrow(a,b]$ 与 $F(b_n)\downarrow F(b)$ 正好对应。
\item $F(x)=\lfloor x\rfloor$（不超过 $x$ 的最大整数）：右连续，单调不减，$F(b)-F(a)=\#(\Z\cap(a,b])$。由唯一性，$\mu_F$ 是 $\Z$ 上的计数测度 $B\mapsto\#(B\cap\Z)$。
\item $F(x)=0$（$x<0$），$F(x)=x$（$0\leqslant x\leqslant1$），$F(x)=1$（$x>1$）：$\mu_F$ 是概率测度，$\mu_F(B)=\lambda(B\cap[0,1])$，称为 $[0,1]$ 上的\term{均匀分布}（uniform distribution）。验证：$B\mapsto\lambda(B\cap[0,1])$ 是测度，在 $(a,b]$ 上的值是 $\lambda((a,b]\cap[0,1])=F(b)-F(a)$（分情况讨论 $a,b$ 与 $0,1$ 的位置，用 (1) 中端点测度为零），再用唯一性。
\item $F(x)=\sum_np_n\one_{[x_n,\infty)}(x)$，$p_n\geqslant0$，$\sum_np_n<\infty$：$F$ 单调不减；右连续性由一致收敛得到（每一项右连续，余项之和 $\sum_{n>N}p_n$ 可以任意小）。$\mu_F=\sum_np_n\delta_{x_n}$，理由同 (1)。
\end{enumerate}
\end{example}

\begin{workedexample}\label{we:extension:atoms}
设 $F$ 是 Stieltjes 测度函数，$\mu_F$ 是对应的 Lebesgue--Stieltjes 测度。记 $F(x-)=\lim_{y\uparrow x}F(y)$。证明：对 $a<b$，
\[
\begin{aligned}
\mu_F(\{x\})&=F(x)-F(x-),&\mu_F([a,b])&=F(b)-F(a-),\\
\mu_F((a,b))&=F(b-)-F(a),&\mu_F([a,b))&=F(b-)-F(a-).
\end{aligned}
\]
并对 $F(x)=x+\one_{[0,\infty)}(x)$ 计算 $\mu_F([-1,0))$，$\mu_F([-1,0])$，$\mu_F(\{0\})$。
\end{workedexample}
\begin{solution}
思路：单点、闭区间、开区间都可以写成左开右闭区间的单调极限，然后用测度的连续性（定理~\ref{thm:measures:basic}(6)(7)）。

\emph{单点。}$(x-\frac1n,x]\downarrow\{x\}$，并且 $\mu_F((x-1,x])=F(x)-F(x-1)<\infty$，由从上连续，
\[
\mu_F(\{x\})=\lim_n\bigl(F(x)-F(x-\tfrac1n)\bigr)=F(x)-F(x-).
\]
最后一步：$F$ 单调，所以左极限存在，并且等于沿任何一个递增趋于 $x$ 的数列的极限。

\emph{闭区间。}$[a,b]=\{a\}+(a,b]$，所以 $\mu_F([a,b])=F(a)-F(a-)+F(b)-F(a)=F(b)-F(a-)$。

\emph{开区间。}$(a,b-\frac1n]\uparrow(a,b)$（$n$ 充分大），由从下连续，$\mu_F((a,b))=\lim_n(F(b-\frac1n)-F(a))=F(b-)-F(a)$。

\emph{左闭右开。}$[a,b)=\{a\}+(a,b)$，所以 $\mu_F([a,b))=F(a)-F(a-)+F(b-)-F(a)=F(b-)-F(a-)$。

\emph{具体计算。}$F(x)=x$（$x<0$），$F(x)=x+1$（$x\geqslant0$）。$F$ 单调不减、右连续，在 $0$ 处有大小为 $1$ 的跳跃：$F(0-)=0$，$F(0)=1$。于是
\[
\begin{aligned}
&\mu_F([-1,0))=F(0-)-F(-1-)=0-(-1)=1,\\
&\mu_F([-1,0])=F(0)-F(-1-)=2,\qquad\mu_F(\{0\})=1.
\end{aligned}
\]
实际上 $\mu_F=\lambda+\delta_0$：右边是测度（\egref{ex:measures:combination}），在 $(a,b]$ 上的值是 $(b-a)+\one_{(a,b]}(0)=F(b)-F(a)$，由唯一性二者相等。

这个例题说明：$F$ 的跳跃对应 $\mu_F$ 在单点上的质量，$F$ 连续的地方单点没有质量。以后讲分布函数时，“$P(X=x)=F(x)-F(x-)$”就是这里的第一个公式。
\end{solution}

\begin{workedexample}\label{we:extension:no-seven}
设 $E$ 为 $[0,1)$ 中十进制展开里不出现数字 $7$ 的数的全体（约定不用以全 $9$ 结尾的展开）。证明 $E$ 是 Borel 集，并且 $\lambda(E)=0$。
\end{workedexample}
\begin{solution}
令 $E_n$ 为前 $n$ 位小数都不是 $7$ 的数的全体。给定前 $n$ 位小数 $d_1,\dots,d_n$，具有这些前 $n$ 位的数恰好组成区间 $[0.d_1\cdots d_n,\ 0.d_1\cdots d_n+10^{-n})$，长度为 $10^{-n}$（这里用到了“不用以全 $9$ 结尾的展开”的约定，它使每个数的展开唯一，并使右端点不属于这个区间）。$E_n$ 是 $9^n$ 个这样的两两不交的区间之并，是 Borel 集，且由有限可加性与例题（第~\pageref{we:extension:atoms} 页）中 $\lambda([a,b))=b-a$，
\[
\lambda(E_n)=9^n\cdot10^{-n}=(0.9)^n .
\]
$E=\bigcap_nE_n$ 是可数交，属于 $\cR$。$E\subseteq E_n$，所以 $\lambda(E)\leqslant(0.9)^n\to0$。

用概率的语言说：在 $[0,1)$ 中均匀地随机取一个数，它的十进制展开中迟早会出现 $7$ 的概率是 $1$。更一般地，每一个有限的数字串都几乎必然在展开中出现；这和第~\ref{lec:independence} 讲的正规数定理是一类问题。这个集合 $E$ 与 Cantor 集是同一类东西：不可数，却是零测集。
\end{solution}

反过来，$\cR$ 上的每个在有界区间上取有限值的测度 $\mu$ 都是 Lebesgue--Stieltjes 测度：令 $F(x)=\mu((0,x])$（$x\geqslant0$），$F(x)=-\mu((x,0])$（$x<0$）。则对 $a<b$，$F(b)-F(a)=\mu((a,b])$（分 $0\leqslant a$、$a<0\leqslant b$、$b<0$ 三种情形验证），于是 $F$ 单调不减；若 $y\downarrow x$，则 $(x,y]\downarrow\varnothing$，且 $\mu((x,x+1])<\infty$，由从上连续，$F(y)-F(x)=\mu((x,y])\to0$，即 $F$ 右连续。由唯一性，$\mu=\mu_F$。两个 Stieltjes 测度函数给出同一个测度，当且仅当它们相差一个常数。

最后，$\R^d$ 上有平行的结论。

\begin{theorem}[$\R^d$ 上的乘积型 Lebesgue--Stieltjes 测度]\label{thm:extension:Rd}
设 $F_1,\dots,F_d$ 是 Stieltjes 测度函数。则 $(\R^d,\cR^d)$ 上存在唯一的测度 $\mu$，使
\[
\mu\bigl((a_1,b_1]\times\cdots\times(a_d,b_d]\bigr)=\prod_{i=1}^d\bigl(F_i(b_i)-F_i(a_i)\bigr)\qquad\text{对一切有限的}\ a_i<b_i
\]
成立。$F_i(x)=x$ 时得到的测度称为 $d$ 维 Lebesgue 测度，记作 $\lambda_d$。
\end{theorem}

证明与定理~\ref{thm:extension:lebesgue-stieltjes} 的证明完全平行，只是有限可加性的验证需要把矩形的有限划分加细为“网格”划分。详细的步骤见作业 B。唯一性部分与一维情形相同：有限矩形连同空集组成生成 $\cR^d$ 的 $\pi$-类（习题~\exref{ex:set-systems:A5}），而 $(-n,n]^d\uparrow\R^d$ 的测度有限。

\topic{$\R^d$ 上一般的分布函数★}

定理~\ref{thm:extension:Rd} 中的测度是“乘积型”的，以后会知道它对应于各坐标独立的情形。一般的 $\R^d$ 上的概率测度 $\mu$ 也由它的 $d$ 维分布函数 $F(x)=\mu((-\infty,x_1]\times\cdots\times(-\infty,x_d])$ 决定（推论~\ref{cor:measures:distribution-function}）。自然的问题是：什么样的 $F:\R^d\to[0,1]$ 是某个概率测度的分布函数？类比一维，我们会猜测下面三条就够了：
\begin{enumerate}
\item[(i)] $F$ 单调不减：若 $x\leqslant y$（即 $x_i\leqslant y_i$ 对一切 $i$），则 $F(x)\leqslant F(y)$；
\item[(ii)] $F$ 右连续：若 $y\downarrow x$（即每个 $y_i\downarrow x_i$），则 $F(y)\to F(x)$；
\item[(iii)] 若 $x$ 的某个坐标趋于 $-\infty$，则 $F(x)\to0$；若 $x$ 的每个坐标都趋于 $+\infty$，则 $F(x)\to1$。
\end{enumerate}
下面的例子说明这是不够的。

\begin{example}[满足 (i)--(iii) 却不是分布函数]\label{ex:extension:bad-F}
在 $\R^2$ 上令
\[
F(x_1,x_2)=\begin{cases}1,&x_1\geqslant1,\ x_2\geqslant1,\\ 2/3,&x_1\geqslant1,\ 0\leqslant x_2<1,\\ 2/3,&x_2\geqslant1,\ 0\leqslant x_1<1,\\ 0,&\text{其他}.\end{cases}
\]
不难逐条验证 (i)--(iii)。假设 $F$ 是某个概率测度 $\mu$ 的分布函数。对矩形 $(a_1,b_1]\times(a_2,b_2]$，由容斥原理，
\[
\mu\bigl((a_1,b_1]\times(a_2,b_2]\bigr)=F(b_1,b_2)-F(a_1,b_2)-F(b_1,a_2)+F(a_1,a_2).
\]
（这是因为 $(a_1,b_1]\times(a_2,b_2]=(-\infty,b]\setminus\bigl[(-\infty,(a_1,b_2)]\cup(-\infty,(b_1,a_2)]\bigr]$，而方括号中两个集合的交是 $(-\infty,a]$。）取 $a_1=a_2=1-\varepsilon$，$b_1=b_2=1$，得
\[
\mu\bigl((1-\varepsilon,1]^2\bigr)=1-\tfrac23-\tfrac23+0=-\tfrac13<0,
\]
这是不可能的。所以 $F$ 不是任何概率测度的分布函数。
\end{example}

这个例子告诉我们，在 $\R^d$ 中，除了单调性，还要求每个矩形按容斥原理算出的“质量”是非负的。对有限矩形 $A=(a_1,b_1]\times\cdots\times(a_d,b_d]$，记 $V=\{a_1,b_1\}\times\cdots\times\{a_d,b_d\}$ 为它的 $2^d$ 个顶点，对顶点 $v$ 令 $\operatorname{sgn}(v)=(-1)^{\#\{i:\ v_i=a_i\}}$，
\[
\Delta_AF=\sum_{v\in V}\operatorname{sgn}(v)F(v).
\]
$d=1$ 时 $\Delta_{(a,b]}F=F(b)-F(a)$，$d=2$ 时就是上面的四项式。第四个条件是
\begin{enumerate}
\item[(iv)] 对一切有限矩形 $A$，$\Delta_AF\geqslant0$。
\end{enumerate}

\begin{theorem}\label{thm:extension:Rd-general}
设 $F:\R^d\to[0,1]$ 满足 (i)--(iv)。则 $(\R^d,\cR^d)$ 上存在唯一的概率测度 $\mu$，使 $\mu(A)=\Delta_AF$ 对一切有限矩形 $A$ 成立。
\end{theorem}
\begin{proof}[证明★]
证明与作业 B（即定理~\ref{thm:extension:Rd} 的证明）完全相同，只有两处需要说明。

\emph{有限可加性。}在作业 B1) 中，乘积的展开换成下面的观察：对函数 $G:\R^d\to\R$，令 $D^{(i)}_{s,t}G(x)=G(x|_{x_i=t})-G(x|_{x_i=s})$（把第 $i$ 个坐标分别换成 $t$、$s$）。则 $\Delta_AF=D^{(1)}_{a_1,b_1}\cdots D^{(d)}_{a_d,b_d}F$，各个 $D^{(i)}$ 都是线性的，作用于不同坐标时可交换。对一维的分点 $a_i=\alpha_0<\cdots<\alpha_n=b_i$，有 $D^{(i)}_{a_i,b_i}=\sum_jD^{(i)}_{\alpha_{j-1},\alpha_j}$（逐项相消）。把每个坐标都这样拆开，就得到网格划分上的可加性。

\emph{连续性。}作业 B3) 中需要 $\Delta_{(a+\delta\mathbf 1,b]}F\to\Delta_{(a,b]}F$（$\delta\downarrow0$）与 $\Delta_{(c,d+\eta\mathbf 1]}F\to\Delta_{(c,d]}F$（$\eta\downarrow0$）。$\Delta$ 是 $2^d$ 个 $F$ 值的有限和，每个顶点的坐标从右边趋于极限顶点，由 (ii) 即得。

单调性（若矩形 $A\subseteq B$，则 $\Delta_AF\leqslant\Delta_BF$）由有限可加性与 (iv) 得到：$B$ 的一个网格划分中含有 $A$。对无界矩形 $A$，令 $\mu(A)=\lim_n\Delta_{A\cap(-n,n]^d}F$（由单调性，极限存在），作业 B 中处理无界矩形的论证照搬过来即可。

这样得到的测度定义在 $\cR^d$ 上，在有限矩形上等于 $\Delta_AF$。由 (iii) 与连续性，$\mu(\R^d)=\lim_n\mu((-n,n]^d)=\lim_n\Delta_{(-n,n]^d}F=1$：展开 $\Delta$，只有所有坐标都取 $n$ 的那个顶点的值趋于 $1$，其余顶点至少有一个坐标为 $-n$，值趋于 $0$。唯一性同定理~\ref{thm:extension:Rd}。
\end{proof}

\begin{example}[乘积型]
若 $F(x)=\prod_{i=1}^dF_i(x_i)$，每个 $F_i$ 是一维的分布函数，则
\[
\Delta_AF=D^{(1)}_{a_1,b_1}\cdots D^{(d)}_{a_d,b_d}\prod_iF_i(x_i)=\prod_{i=1}^d\bigl(F_i(b_i)-F_i(a_i)\bigr)\geqslant0,
\]
因为 $D^{(i)}$ 只作用在第 $i$ 个因子上。所以 (iv) 自动成立，这时定理~\ref{thm:extension:Rd-general} 就是定理~\ref{thm:extension:Rd} 的概率测度情形。$F_i(x)=x$ 的情形（不是概率测度，但论证相同）给出 $d$ 维 Lebesgue 测度。
\end{example}

\topic{本讲小结}

\begin{enumerate}
\item 构造测度的标准路线：在半代数 $\cS$ 上给出集函数 $\mu$，验证有限可加性 (i) 与可数次可加性 (ii)；它自动延拓为代数 $\overline\cS$ 上的测度（定理~\ref{thm:extension:semialgebra}），再由 Carathéodory 延拓定理延拓到 $\sigma(\cS)$ 上（定理~\ref{thm:extension:main}、推论~\ref{cor:extension:semialgebra}）。$\sigma$-有限时延拓唯一。
\item 延拓的方法是外测度：用可数个代数中的集合从外面盖住，取测度之和的下确界。“可数”是 Lebesgue 与 Jordan、Riemann 的区别所在。
\item 对单调不减、右连续的 $F$，存在唯一的测度 $\mu_F$ 使 $\mu_F((a,b])=F(b)-F(a)$（定理~\ref{thm:extension:lebesgue-stieltjes}）。$F(x)=x$ 给出 Lebesgue 测度；$F$ 的跳跃对应单点的质量（例题，第~\pageref{we:extension:atoms} 页）。证明中唯一的分析是：紧性加右连续性给出可数次可加性。
\item 下一讲以后，我们总是取分布函数为 $F$ 的 $\mu_F$ 作为随机变量的分布（第~\ref{lec:distributions} 讲）；$\R^d$ 上要求矩形的容斥“质量”非负（定理~\ref{thm:extension:Rd-general}）。
\end{enumerate}

% =====================================================================
\homework{$d$ 维 Lebesgue 测度，可测集的逼近，不可测集}
\label{hw:extension}

\exercisepart{基本习题}

\begin{exercises}{A}{课堂内容的补充}
\item\label{ex:extension:A1} 设 $\cS$ 是 $\Omega_1$ 上的半代数，$\cT$ 是 $\Omega_2$ 上的半代数。证明 $\{S\times T: S\in\cS,\ T\in\cT\}$ 是 $\Omega_1\times\Omega_2$ 上的半代数。由此重新证明 $\cS_d$ 是半代数。
\begin{solution}
记这个族为 $\cS\times\cT$。

\emph{交集。}$(S\times T)\cap(S'\times T')=(S\cap S')\times(T\cap T')$：$(x,y)$ 属于左边当且仅当 $x\in S$，$x\in S'$，$y\in T$，$y\in T'$。$S\cap S'\in\cS$，$T\cap T'\in\cT$。

\emph{补集。}$(x,y)\notin S\times T$ 当且仅当 $x\notin S$，或者 $x\in S$ 而 $y\notin T$。所以
\[
(S\times T)^c=(S^c\times\Omega_2)+(S\times T^c),
\]
两部分不交。写 $S^c=\sum_iS_i$，$T^c=\sum_jT_j$（有限不交并，$S_i\in\cS$，$T_j\in\cT$），又 $\Omega_2=T+\sum_jT_j$。于是
\[
(S\times T)^c=\sum_i(S_i\times T)+\sum_{i,j}(S_i\times T_j)+\sum_j(S\times T_j),
\]
每一项都属于 $\cS\times\cT$，并且两两不交（第一坐标或第二坐标落在不交的集合中）。

\emph{$\cS_d$。}把 $\R^d$ 看成 $\R^{d-1}\times\R$，则 $\cS_d=\cS_{d-1}\times\cS_1$（某个因子为空集时乘积为空集）。由 $\cS_1$ 是半代数（\egref{ex:extension:S1}）与归纳法即得。
\end{solution}

\item\label{ex:extension:A2} 设 $\mu^*$ 是外测度。
\begin{questions}
\item 证明：若对一切 $\mu^*(F)<\infty$ 的 $F$ 都有 $\mu^*(F)\geqslant\mu^*(F\cap E)+\mu^*(F\cap E^c)$，则 $E$ 是 $\mu^*$-可测的。
\begin{solution}
要证对一切 $F$ 有 $\mu^*(F)=\mu^*(F\cap E)+\mu^*(F\cap E^c)$。由次可加性（$F=(F\cap E)\cup(F\cap E^c)$），“$\leqslant$”总成立。“$\geqslant$”：若 $\mu^*(F)<\infty$，这是假设；若 $\mu^*(F)=\infty$，左边是 $\infty$，不等式自动成立。
\end{solution}
\item 证明：若 $E_1\subseteq E_2\subseteq\cdots$ 都是 $\mu^*$-可测的，则 $\mu^*(G\cap\bigcup_nE_n)=\lim_n\mu^*(G\cap E_n)$ 对一切 $G$ 成立。
\begin{solution}
令 $E=\bigcup_nE_n$，$D_1=E_1$，$D_n=E_n\setminus E_{n-1}$。由定理~\ref{thm:extension:caratheodory} 的证明 (b)，$D_n\in\cM$，它们两两不交，$\bigcup_{i\leqslant n}D_i=E_n$，$\bigcup_iD_i=E$。由证明中的 (c)（对试验集 $G$），
\[
\mu^*(G\cap E_n)=\sum_{i\leqslant n}\mu^*(G\cap D_i).
\]
在 \eqref{eq:extension:countable} 中把 $G$ 换成 $G\cap E$，并注意 $G\cap E\cap D_i=G\cap D_i$，$G\cap E\cap E^c=\varnothing$：
\[
\mu^*(G\cap E)\geqslant\sum_{i=1}^\infty\mu^*(G\cap D_i)=\lim_n\mu^*(G\cap E_n).
\]
反向：$G\cap E_n\subseteq G\cap E$，由单调性 $\mu^*(G\cap E_n)\leqslant\mu^*(G\cap E)$。两者合起来得到等式。

注意 $G$ 是任意集合，不必可测。这个结论说明外测度对“被可测集截下的部分”有从下连续性。
\end{solution}
\end{questions}

\item\label{ex:extension:A3} 证明 Lebesgue 测度的\term{伸缩性}：对 $A\in\cR$ 与 $c\neq0$，$cA=\{ca:a\in A\}\in\cR$，且 $\lambda(cA)=\abs c\lambda(A)$。（提示：分 $c>0$ 与 $c=-1$ 两种情形，注意 $-(a,b]=[-b,-a)$。）
\begin{solution}
$cA\in\cR$ 已在第~\ref{lec:set-systems} 讲的例题（第~\pageref{we:set-systems:translation} 页）中证明。

\emph{$c>0$。}令 $\nu(A)=\frac1c\lambda(cA)$，$A\in\cR$。$x\mapsto cx$ 是双射，所以不交的集合映成不交的集合，$c\bigcup_nA_n=\bigcup_n cA_n$，于是 $\nu$ 是测度。$\nu((a,b])=\frac1c\lambda((ca,cb])=\frac1c(cb-ca)=b-a$。由定理~\ref{thm:extension:lebesgue-stieltjes} 的唯一性，$\nu=\lambda$，即 $\lambda(cA)=c\lambda(A)$。

\emph{$c=-1$。}同理 $\nu(A)=\lambda(-A)$ 是测度，$\nu((a,b])=\lambda([-b,-a))=b-a$（例题，第~\pageref{we:extension:atoms} 页），所以 $\nu=\lambda$。

\emph{$c<0$。}$cA=-(\abs cA)$，由前两种情形 $\lambda(cA)=\lambda(\abs cA)=\abs c\lambda(A)$。
\end{solution}

\item\label{ex:extension:A4} \term{Cantor 集}（Cantor set）：令 $C_0=[0,1]$，$C_{n+1}$ 由 $C_n$ 的每个闭区间去掉中间的开三分之一得到，$C=\bigcap_nC_n$。
\begin{questions}
\item 证明 $C$ 是闭集，$C\in\cR$，且 $\lambda(C)=0$。
\begin{solution}
对 $n$ 归纳可知，$C_n$ 是 $2^n$ 个两两不交、长度为 $3^{-n}$ 的闭区间的并：每个长为 $3^{-n}$ 的闭区间去掉中间的开三分之一，剩下两个长为 $3^{-n-1}$ 的不交闭区间。所以 $C_n$ 是闭集，$\lambda(C_n)=2^n3^{-n}=(\frac23)^n$（闭区间的测度等于长度，例题第~\pageref{we:extension:atoms} 页）。$C$ 是闭集之交，是闭集，从而是 Borel 集。$C\subseteq C_n$，$\lambda(C)\leqslant(\frac23)^n\to0$。
\end{solution}
\item 证明 $x\in C$ 当且仅当 $x$ 有一个只用数字 $0,2$ 的三进制展开 $x=\sum_{n\geqslant1}\varepsilon_n3^{-n}$，$\varepsilon_n\in\{0,2\}$。
\begin{solution}
\emph{第一步：$C_n$ 的描述。}对 $n$ 归纳证明
\[
C_n=\bigcup_{\varepsilon_1,\dots,\varepsilon_n\in\{0,2\}}\bigl[c(\varepsilon),\,c(\varepsilon)+3^{-n}\bigr],\qquad c(\varepsilon)=\sum_{k=1}^n\varepsilon_k3^{-k}.
\]
$n=0$ 时是 $[0,1]$。区间 $[c,c+3^{-n}]$ 去掉中间三分之一后剩下 $[c,c+3^{-n-1}]$ 与 $[c+2\cdot3^{-n-1},c+3^{-n}]$，分别对应 $\varepsilon_{n+1}=0$ 与 $\varepsilon_{n+1}=2$。

\emph{第二步：“$\Leftarrow$”。}设 $x=\sum_k\varepsilon_k3^{-k}$，$\varepsilon_k\in\{0,2\}$。对每个 $n$，记 $c_n=\sum_{k\leqslant n}\varepsilon_k3^{-k}$，则 $0\leqslant x-c_n\leqslant\sum_{k>n}2\cdot3^{-k}=3^{-n}$，所以 $x\in[c_n,c_n+3^{-n}]\subseteq C_n$。这对一切 $n$ 成立，$x\in C$。

\emph{第三步：“$\Rightarrow$”。}设 $x\in C$。对每个 $n$，$C_n$ 的 $2^n$ 个区间两两不交，$x$ 恰好属于其中一个，记其标号为 $(\varepsilon_1^{(n)},\dots,\varepsilon_n^{(n)})$。第 $n+1$ 层中含 $x$ 的区间包含于第 $n$ 层中含 $x$ 的区间，而第一步说明前者的标号是后者的标号加上一位，所以 $\varepsilon_k^{(n)}$ 不依赖于 $n\geqslant k$，记为 $\varepsilon_k$。于是 $x\in[c_n,c_n+3^{-n}]$ 对一切 $n$ 成立，$c_n\to\sum_k\varepsilon_k3^{-k}$，从而 $x=\sum_k\varepsilon_k3^{-k}$。
\end{solution}
\item 由 A4-2) 证明 $C$ 与 $[0,1]$ 等势，从而是不可数的零测集。
\begin{solution}
令 $\Phi(\varepsilon)=\sum_k\varepsilon_k3^{-k}$，$\varepsilon\in\{0,2\}^\N$。由 A4-2)，$\Phi$ 是到 $C$ 上的满射。它是单射：若 $\varepsilon\neq\varepsilon'$，设第一个不同的位置是 $k$，比如 $\varepsilon_k=2$，$\varepsilon'_k=0$，则
\[
\Phi(\varepsilon)-\Phi(\varepsilon')\geqslant2\cdot3^{-k}-\sum_{j>k}2\cdot3^{-j}=2\cdot3^{-k}-3^{-k}=3^{-k}>0.
\]
所以 $\abs C=\abs{\{0,2\}^\N}=\abs{\{0,1\}^\N}=\abs{[0,1]}$（最后一个等式是二进制展开给出的标准事实）。$C$ 不可数，由 A4-1) 是零测集。

这和“可数集测度为零”对照：测度为零的集合可以很大（与 $\R$ 等势），测度衡量的不是点的多少。
\end{solution}
\end{questions}

\item\label{ex:extension:A5} 设 $\lambda^*$ 是长度在代数 $\overline{\cS_1}$ 上诱导的外测度。
\begin{questions}
\item 用外测度的定义直接证明：$\lambda^*(\Q\cap[0,1])=0$。
\begin{solution}
把 $\Q\cap[0,1]$ 排成 $q_1,q_2,\dots$。任取 $\varepsilon>0$，令 $A_n=(q_n-\varepsilon2^{-n-1},\,q_n+\varepsilon2^{-n-1}]\in\cS_1$。$\{A_n\}$ 覆盖 $\Q\cap[0,1]$，长度之和为 $\sum_n\varepsilon2^{-n}=\varepsilon$。由外测度的定义 $\lambda^*(\Q\cap[0,1])\leqslant\varepsilon$。$\varepsilon$ 任意，所以等于 $0$。
\end{solution}
\item 证明对任意 $E\subseteq\R$，$\lambda^*(E)=\inf\{\lambda(U): U\supseteq E,\ U\ \text{开}\}$。
\begin{solution}
\emph{“$\leqslant$”。}若 $U\supseteq E$ 是开集，则 $U\in\cR\subseteq\cM(\lambda^*)$，$\lambda(U)=\lambda^*(U)\geqslant\lambda^*(E)$（定理~\ref{thm:extension:main} 的证明中 $\lambda=\lambda^*$ 在 $\cR$ 上；外测度单调）。

\emph{“$\geqslant$”。}若 $\lambda^*(E)=\infty$，显然。否则任取 $\varepsilon>0$，由定义取 $A_n\in\overline{\cS_1}$ 覆盖 $E$，$\sum_n\lambda(A_n)\leqslant\lambda^*(E)+\varepsilon<\infty$。每个 $A_n$ 是有限个左开右闭区间的不交并，这些区间都有界（否则 $\lambda(A_n)=\infty$）。把全体这样的区间排成 $(a_k,b_k]$，$k\in\N$，则它们覆盖 $E$，$\sum_k(b_k-a_k)\leqslant\lambda^*(E)+\varepsilon$。令 $U=\bigcup_k(a_k,\,b_k+\varepsilon2^{-k})$，这是开集，包含 $E$，并且由次可加性
\[
\lambda(U)\leqslant\sum_k\bigl(b_k-a_k+\varepsilon2^{-k}\bigr)\leqslant\lambda^*(E)+2\varepsilon .
\]
$\varepsilon$ 任意，得“$\geqslant$”。
\end{solution}
\end{questions}

\item\label{ex:extension:A6} 证明\egref{ex:extension:LS-examples} 中的 (3)：分情况验证 $\lambda((a,b]\cap[0,1])=F(b)-F(a)$。
\begin{solution}
与其逐一讨论 $a,b$ 与 $0,1$ 的相对位置，不如先算 $(-\infty,x]\cap[0,1]$ 的测度：$x<0$ 时它是空集，测度 $0$；$0\leqslant x\leqslant1$ 时它是 $[0,x]$，测度 $x$；$x>1$ 时它是 $[0,1]$，测度 $1$。所以 $\lambda((-\infty,x]\cap[0,1])=F(x)$ 对一切 $x$ 成立。而
\[
(a,b]\cap[0,1]=\bigl((-\infty,b]\cap[0,1]\bigr)\setminus\bigl((-\infty,a]\cap[0,1]\bigr),
\]
后者包含于前者且测度有限，由定理~\ref{thm:measures:basic}(3)，$\lambda((a,b]\cap[0,1])=F(b)-F(a)$。这正是六种情形的统一写法。
\end{solution}

\item\label{ex:extension:A7} 设 $F,G$ 是 Stieltjes 测度函数。证明 $\mu_F(B)\leqslant\mu_G(B)$ 对一切 $B\in\cR$ 成立，当且仅当 $G-F$ 单调不减。
\begin{solution}
“$\Rightarrow$”：对 $x<y$ 取 $B=(x,y]$，$F(y)-F(x)\leqslant G(y)-G(x)$，即 $(G-F)(x)\leqslant(G-F)(y)$。

“$\Leftarrow$”：$H=G-F$ 单调不减，又是两个右连续函数之差，右连续，所以是 Stieltjes 测度函数。$\mu_F+\mu_H$ 是测度，在 $(a,b]$ 上取值 $F(b)-F(a)+H(b)-H(a)=G(b)-G(a)$，由唯一性 $\mu_G=\mu_F+\mu_H$，于是 $\mu_G(B)=\mu_F(B)+\mu_H(B)\geqslant\mu_F(B)$。
\end{solution}

\item\label{ex:extension:A8} 设 $\cA$ 是 $\N$ 上有限集与余有限集组成的代数，$\mu$ 如\egref{ex:measures:finitely-additive}（有限集上为 $0$，余有限集上为 $1$）。证明 $\mu$ 不能延拓为 $2^\N$（甚至 $\sigma(\cA)$）上的测度。它违反了定理~\ref{thm:extension:main} 的哪个条件？
\begin{solution}
若 $\nu$ 是 $\sigma(\cA)$ 上的测度且在 $\cA$ 上等于 $\mu$，则 $1=\mu(\N)=\nu(\N)=\sum_n\nu(\{n\})=\sum_n\mu(\{n\})=0$，矛盾。$2^\N\supseteq\sigma(\cA)$ 上更不可能。

它违反的是定理的前提：$\mu$ 不是代数 $\cA$ 上的测度。$\{n\}\in\cA$ 两两不交，并 $\N\in\cA$，但 $\mu(\N)\neq\sum_n\mu(\{n\})$。延拓定理的唯一前提就是代数上的可数可加性，这个例子说明它是必需的。
\end{solution}

\item\label{ex:extension:A9} 设 $\mu$ 是 $(\R,\cR)$ 上的测度，对一切 $x$ 有 $\mu(\{x\})=0$，且 $\mu$ 在有界区间上有限。证明对应的 Stieltjes 测度函数 $F$ 是连续的。
\begin{solution}
$F$ 右连续，只需证左连续。由例题（第~\pageref{we:extension:atoms} 页），$F(x)-F(x-)=\mu(\{x\})=0$。（那里的证明用到 $(x-\frac1n,x]\downarrow\{x\}$ 与从上连续，有限性由 $\mu$ 在有界区间上有限保证。）
\end{solution}

\item\label{ex:extension:A10} 设 $\mu$ 是 $(\R,\cR)$ 上的测度，满足 $\mu(A+x)=\mu(A)$ 对一切 $A\in\cR$，$x\in\R$ 成立，并且 $\mu((0,1])=1$。证明 $\mu=\lambda$。（提示：先对 $(0,1/n]$，再对端点为有理数的区间，最后对一切区间计算 $\mu$。）
\begin{solution}
\emph{第一步。}$(0,1]=\sum_{k=1}^n(\frac{k-1}n,\frac kn]$，每一项是 $(0,\frac1n]$ 平移 $\frac{k-1}n$ 所得，测度相同。由有限可加性 $1=n\mu((0,\frac1n])$，$\mu((0,\frac1n])=\frac1n$。

\emph{第二步：有理端点。}设 $a<b$ 有理，写成同分母 $a=\frac pn$，$b=\frac qn$。$(a,b]$ 是 $q-p$ 个 $(0,\frac1n]$ 的平移的不交并，$\mu((a,b])=\frac{q-p}n=b-a$。

\emph{第三步：一般区间。}设 $a<b$ 为实数。取有理数 $a_m\downarrow a$，$a_m<b$，以及有理数 $b_k\downarrow b$。对固定的 $m$，$(a_m,b_k]\downarrow(a_m,b]$，且 $\mu((a_m,b_1])<\infty$，由从上连续 $\mu((a_m,b])=\lim_k(b_k-a_m)=b-a_m$。再由 $(a_m,b]\uparrow(a,b]$ 与从下连续，$\mu((a,b])=\lim_m(b-a_m)=b-a$。

\emph{第四步。}$\mu$ 与 $\lambda$ 在一切 $(a,b]$ 上相等，由第~\ref{lec:measures} 讲的例题（第~\pageref{we:measures:uniqueness-lebesgue} 页），$\mu=\lambda$。
\end{solution}
\end{exercises}

\begin{probleme}{B}{$d$ 维 Lebesgue--Stieltjes 测度的构造}
本题证明定理~\ref{thm:extension:Rd}。设 $F_1,\dots,F_d$ 是 Stieltjes 测度函数，对矩形 $R=\prod_i(a_i,b_i]$（$-\infty\leqslant a_i<b_i\leqslant\infty$）令 $\mu(R)=\prod_i(F_i(b_i)-F_i(a_i))$，约定 $0\cdot\infty=0$，并令 $\mu(\varnothing)=0$。
\begin{questions}
\item 称 $R=\sum_kR_k$ 是\term{网格划分}（regular subdivision），若对每个 $i$ 存在 $a_i=\alpha_{i,0}<\alpha_{i,1}<\cdots<\alpha_{i,n_i}=b_i$，使 $\{R_k\}$ 恰好是全体 $\prod_i(\alpha_{i,j_i-1},\alpha_{i,j_i}]$（$1\leqslant j_i\leqslant n_i$）。证明对网格划分有 $\mu(R)=\sum_k\mu(R_k)$。（提示：把乘积 $\prod_i\sum_{j}(F_i(\alpha_{i,j})-F_i(\alpha_{i,j-1}))$ 展开。）
\begin{solution}
对每个 $i$，由逐项相消，
\[
F_i(b_i)-F_i(a_i)=\sum_{j=1}^{n_i}\bigl(F_i(\alpha_{i,j})-F_i(\alpha_{i,j-1})\bigr).
\]
（端点为 $\pm\infty$、$F_i(\pm\infty)$ 为无穷时，两边都是 $\infty$：相应的首项或末项为 $\infty$，其余项非负。）把 $d$ 个这样的和相乘并展开（$[0,\infty]$ 中约定 $0\cdot\infty=0$ 后分配律仍成立）：
\[
\mu(R)=\prod_{i=1}^d\sum_{j=1}^{n_i}\Delta_{i,j}=\sum_{j_1,\dots,j_d}\prod_{i=1}^d\Delta_{i,j_i},\qquad\Delta_{i,j}=F_i(\alpha_{i,j})-F_i(\alpha_{i,j-1}).
\]
右边每一项恰好是网格中一个小矩形 $\prod_i(\alpha_{i,j_i-1},\alpha_{i,j_i}]$ 的 $\mu$ 值，所以右边等于 $\sum_k\mu(R_k)$。
\end{solution}
\item 设 $R=\sum_{k=1}^mR_k$ 是任意有限划分。取每个坐标上所有 $R_k$ 的端点，得到 $R$ 的一个网格划分 $\{Q_l\}$。证明每个 $R_k$ 是某些 $Q_l$ 的不交并，并且这些 $Q_l$ 构成 $R_k$ 的网格划分。由 B1) 推出条件 (i)。
\begin{solution}
对每个 $i$，把 $a_i,b_i$ 以及所有 $R_k$ 的第 $i$ 个坐标区间的端点从小到大排成 $\alpha_{i,0}<\cdots<\alpha_{i,n_i}$，得到网格 $\{Q_l\}$。

\emph{每个小格要么含于 $R_k$，要么与 $R_k$ 不交。}看第 $i$ 个坐标：小格的区间 $(\alpha_{i,j-1},\alpha_{i,j}]$ 与 $R_k$ 的区间 $(c_i,d_i]$ 的端点都是网格点，所以前者或者含于后者，或者与后者不交。若所有坐标都“含于”，则 $Q_l\subseteq R_k$；只要有一个坐标“不交”，则 $Q_l\cap R_k=\varnothing$。

由于 $R_k$ 两两不交且并为 $R$，而小格非空、含于 $R$，每个小格恰好含于一个 $R_k$。所以 $R_k$ 是含于它的小格的不交并。含于 $R_k$ 的小格在每个坐标上对应 $(c_i,d_i]$ 中的网格点 $c_i=\alpha_{i,s}<\cdots<\alpha_{i,t}=d_i$，因此构成 $R_k$ 的网格划分。

由 B1) 用两次：
\[
\mu(R)=\sum_l\mu(Q_l)=\sum_k\sum_{l:\,Q_l\subseteq R_k}\mu(Q_l)=\sum_k\mu(R_k).
\]
这就是条件 (i)。
\end{solution}
\item 设 $a_i,b_i$ 都有限，$R=(a,b]$ 被可数个矩形 $(a^{(k)},b^{(k)}]$ 覆盖（这里 $(x,y]$ 表示 $\prod_i(x_i,y_i]$）。仿照定理~\ref{thm:extension:lebesgue-stieltjes} 的证明，取 $\delta>0$ 与 $\eta_k>0$，使 $\mu((a+\delta\mathbf 1,b])>\mu((a,b])-\varepsilon$，$\mu((a^{(k)},b^{(k)}+\eta_k\mathbf 1])<\mu((a^{(k)},b^{(k)}])+\varepsilon2^{-k}$，其中 $\mathbf 1=(1,\dots,1)$。说明这样的 $\delta,\eta_k$ 存在。
\begin{solution}
先像一维那样，把每个 $(a^{(k)},b^{(k)}]$ 换成它与 $R$ 的交（仍是矩形或空集，空集去掉）。这不改变覆盖关系，并且 $\mu$ 值不增（每个坐标上的区间变小，$F_i$ 的增量变小）。所以不妨设覆盖矩形都含于 $R$，从而都有限。

$\mu((a+\delta\mathbf 1,b])=\prod_i(F_i(b_i)-F_i(a_i+\delta))$。$\delta\downarrow0$ 时，由 $F_i$ 右连续，每个因子趋于 $F_i(b_i)-F_i(a_i)$，有限个有限数的乘积连续，所以乘积趋于 $\mu((a,b])$，$\delta$ 充分小时满足要求。同理 $\mu((a^{(k)},b^{(k)}+\eta\mathbf 1])=\prod_i(F_i(b_i^{(k)}+\eta)-F_i(a_i^{(k)}))\to\mu((a^{(k)},b^{(k)}])$（$\eta\downarrow0$），取 $\eta_k$ 充分小即可。
\end{solution}
\item 用 $[a+\delta\mathbf 1,b]$ 的紧性与引理~\ref{lem:extension:finite-additive}(3) 完成条件 (ii) 在有限矩形情形的证明；再像一维那样处理无界矩形。
\begin{solution}
\emph{有限矩形。}$[a+\delta\mathbf 1,b]\subseteq(a,b]\subseteq\bigcup_k(a^{(k)},b^{(k)}]\subseteq\bigcup_k(a^{(k)},b^{(k)}+\eta_k\mathbf 1)$，右边是开矩形。$[a+\delta\mathbf 1,b]$ 是 $\R^d$ 中的有界闭集，由 Heine--Borel 定理，有有限个 $k_1,\dots,k_J$ 使
\[
(a+\delta\mathbf 1,b]\subseteq[a+\delta\mathbf 1,b]\subseteq\bigcup_{j\leqslant J}(a^{(k_j)},b^{(k_j)}+\eta_{k_j}\mathbf 1].
\]
B2) 给出了条件 (i)，所以引理~\ref{lem:extension:finite-additive}(3) 可用：
\[
\mu((a,b])-\varepsilon<\mu((a+\delta\mathbf 1,b])\leqslant\sum_{j\leqslant J}\mu\bigl((a^{(k_j)},b^{(k_j)}+\eta_{k_j}\mathbf 1]\bigr)\leqslant\sum_k\bigl(\mu((a^{(k)},b^{(k)}])+\varepsilon2^{-k}\bigr).
\]
于是 $\mu(R)<\sum_k\mu(R_k)+2\varepsilon$，$\varepsilon$ 任意，得 $\mu(R)\leqslant\sum_k\mu(R_k)$。这比条件 (ii) 更强：覆盖不必不交。

\emph{无界矩形。}若 $R$ 的某个因子 $F_i(b_i)-F_i(a_i)=0$，则 $\mu(R)=0$，无需证明。否则对任意有限矩形 $R'=(A,B]\subseteq R$，它也被 $\{R_k\}$ 覆盖，由上一段 $\mu(R')\leqslant\sum_k\mu(R_k)$。令 $A_i\downarrow a_i$，$B_i\uparrow b_i$（有限的端点保持不动），每个因子 $F_i(B_i)-F_i(A_i)$ 单调趋于 $F_i(b_i)-F_i(a_i)\in(0,\infty]$，乘积趋于 $\mu(R)$（各因子的极限都是正数，不出现 $0\cdot\infty$）。所以 $\mu(R)\leqslant\sum_k\mu(R_k)$。
\end{solution}
\item 由推论~\ref{cor:extension:semialgebra} 与习题~\exref{ex:set-systems:A5} 完成定理的证明。
\begin{solution}
B2)、B4) 验证了条件 (i)(ii)。$(-n,n]^d\in\cS_d$，$\mu((-n,n]^d)=\prod_i(F_i(n)-F_i(-n))<\infty$，并为 $\R^d$。由推论~\ref{cor:extension:semialgebra}，存在唯一的 $\sigma(\cS_d)$ 上的测度在 $\cS_d$ 上等于 $\mu$。由习题~\exref{ex:set-systems:A5}，$\sigma(\cS_d)=\cR^d$。至于定理中只要求在有限矩形上相等的唯一性：有限矩形连同空集是生成 $\cR^d$ 的 $\pi$-类，含 $(-n,n]^d$，由定理~\ref{thm:measures:uniqueness} 即得。
\end{solution}
\item 证明 $\lambda_d$ 是平移不变的，并且 $\lambda_d(\{x\})=0$，$\lambda_d(\text{超平面}\ \{x: x_1=c\})=0$。
\begin{solution}
\emph{平移不变。}与第~\ref{lec:set-systems} 讲的例题（第~\pageref{we:set-systems:translation} 页）相同，$A\in\cR^d$ 时 $A+x\in\cR^d$。$\nu(A)=\lambda_d(A+x)$ 是测度，在矩形上 $\nu(\prod_i(a_i,b_i])=\lambda_d(\prod_i(a_i+x_i,b_i+x_i])=\prod_i(b_i-a_i)$，由唯一性 $\nu=\lambda_d$。

\emph{单点。}$\{x\}\subseteq(x-\frac1n\mathbf 1,x]$，$\lambda_d(\{x\})\leqslant n^{-d}\to0$。

\emph{超平面。}$H=\{x:x_1=c\}=\bigcup_m H_m$，$H_m=\{c\}\times(-m,m]^{d-1}$。$H_m\subseteq(c-\frac1n,c]\times(-m,m]^{d-1}$，$\lambda_d(H_m)\leqslant\frac1n(2m)^{d-1}\to0$（$n\to\infty$）。可数个零测集之并是零测集。
\end{solution}
\end{questions}
\end{probleme}

\begin{probleme}{C}{可测集的逼近与完备化}
设 $\mu$ 是代数 $\cA$ 上的 $\sigma$-有限测度，$\mu^*$ 是诱导的外测度，$\cM=\cM(\mu^*)$。记 $\cA_\sigma$ 为 $\cA$ 中集合的可数并的全体，$\cA_{\sigma\delta}$ 为 $\cA_\sigma$ 中集合的可数交的全体。
\begin{questions}
\item 设 $\mu^*(E)<\infty$，$\varepsilon>0$。证明存在 $A\in\cA_\sigma$，$A\supseteq E$，$\mu^*(A)\leqslant\mu^*(E)+\varepsilon$。
\begin{solution}
由外测度的定义，取 $A_n\in\cA$ 覆盖 $E$，$\sum_n\mu(A_n)\leqslant\mu^*(E)+\varepsilon$。令 $A=\bigcup_nA_n\in\cA_\sigma$。由次可加性与 $\mu^*=\mu$ 在 $\cA$ 上，$\mu^*(A)\leqslant\sum_n\mu^*(A_n)=\sum_n\mu(A_n)\leqslant\mu^*(E)+\varepsilon$。
\end{solution}
\item 在 C1) 中进一步设 $E\in\cM$。证明存在 $B\in\cA$ 使 $\mu^*(E\triangle B)\leqslant2\varepsilon$。（提示：$A=\bigcup_iA_i$，取有限并 $B=\bigcup_{i\leqslant n}A_i$。）
\begin{solution}
$A\in\sigma(\cA)\subseteq\cM$，$E\in\cM$，$\mu^*$ 在 $\cM$ 上是测度，且 $\mu^*(E)<\infty$，所以 $\mu^*(A\setminus E)=\mu^*(A)-\mu^*(E)\leqslant\varepsilon$。令 $B_n=\bigcup_{i\leqslant n}A_i\in\cA$，$B_n\uparrow A$，$\mu^*(A)<\infty$，由从上连续 $\mu^*(A\setminus B_n)\to0$，取 $n$ 使它 $\leqslant\varepsilon$，令 $B=B_n$。由于 $B\subseteq A$，$E\subseteq A$：
\[
B\setminus E\subseteq A\setminus E,\qquad E\setminus B\subseteq A\setminus B,
\]
所以 $\mu^*(E\triangle B)\leqslant\varepsilon+\varepsilon=2\varepsilon$。
\end{solution}
\item 设 $\mu^*(E)<\infty$。证明存在 $G\in\cA_{\sigma\delta}$，$G\supseteq E$，$\mu^*(G)=\mu^*(E)$。
\begin{solution}
由 C1)，对每个 $k$ 取 $A_k\in\cA_\sigma$，$A_k\supseteq E$，$\mu^*(A_k)\leqslant\mu^*(E)+\frac1k$。令 $G=\bigcap_kA_k\in\cA_{\sigma\delta}$。$E\subseteq G\subseteq A_k$，由单调性 $\mu^*(E)\leqslant\mu^*(G)\leqslant\mu^*(E)+\frac1k$ 对一切 $k$ 成立。
\end{solution}
\item 设 $E\in\cM$，$\mu^*(E)<\infty$，$G$ 如 C3)。证明 $\mu^*(G\setminus E)=0$。
\begin{solution}
$G,E\in\cM$，$E\subseteq G$，$\mu^*(E)<\infty$，由测度的减法 $\mu^*(G\setminus E)=\mu^*(G)-\mu^*(E)=0$。

注意 $E\notin\cM$ 时结论可以不成立：此时 $\mu^*(G)=\mu^*(E)$ 不能推出差集的外测度为零，因为 $\mu^*$ 在 $2^\Omega$ 上不是可加的。
\end{solution}
\item 利用 $\sigma$-有限性去掉 C4) 中 $\mu^*(E)<\infty$ 的假设：证明 $E\in\cM$ 当且仅当存在 $G\in\cA_{\sigma\delta}\subseteq\sigma(\cA)$ 与 $\mu^*(N)=0$ 的集合 $N$，使 $E=G\setminus N$。
\begin{solution}
“$\Leftarrow$”：$G\in\sigma(\cA)\subseteq\cM$；$\mu^*(N)=0$，由 Carathéodory 定理 $N\in\cM$；所以 $E=G\setminus N\in\cM$。

“$\Rightarrow$”：取 $\Omega_m\in\cA$ 两两不交，$\bigcup_m\Omega_m=\Omega$，$\mu(\Omega_m)<\infty$（$\sigma$-有限性，见第~\ref{lec:measures} 讲定义~\ref{def:measures:finite} 之后的说明）。令 $E_m=E\cap\Omega_m\in\cM$，$\mu^*(E_m)<\infty$。对每个 $m,k$，由 C1) 取 $A_{m,k}\in\cA_\sigma$，$A_{m,k}\supseteq E_m$，$\mu^*(A_{m,k})\leqslant\mu^*(E_m)+\frac{2^{-m}}k$；像 C4) 一样，$\mu^*(A_{m,k}\setminus E_m)\leqslant\frac{2^{-m}}k$。令
\[
A_k=\bigcup_mA_{m,k}\in\cA_\sigma,\qquad G=\bigcap_kA_k\in\cA_{\sigma\delta}.
\]
$E\subseteq A_k$，并且 $A_k\setminus E\subseteq\bigcup_m(A_{m,k}\setminus E_m)$（若 $x\in A_{m,k}$ 且 $x\notin E$，则 $x\notin E_m$），所以 $\mu^*(A_k\setminus E)\leqslant\sum_m\frac{2^{-m}}k=\frac1k$。$G\setminus E\subseteq A_k\setminus E$ 对一切 $k$ 成立，$\mu^*(G\setminus E)=0$。令 $N=G\setminus E$，则 $E=G\setminus N$。

（这里没有直接对每个 $E_m$ 用 C3)，因为 $\cA_{\sigma\delta}$ 中集合的可数并不一定还在 $\cA_{\sigma\delta}$ 中；先在 $\cA_\sigma$ 这一层取并就避开了这个问题。）
\end{solution}
\item 由 C5) 与上一讲作业 B 证明：$(\Omega,\cM,\mu^*)$ 是 $(\Omega,\sigma(\cA),\mu^*|_{\sigma(\cA)})$ 的完备化。
\begin{solution}
记 $\overline{\sigma(\cA)}$ 为完备化后的 $\sigma$-代数，$\cN$ 为可略集（含于 $\sigma(\cA)$ 中零测集的集合）。

\emph{$\overline{\sigma(\cA)}\subseteq\cM$。}$\sigma(\cA)\subseteq\cM$；若 $B\in\cN$，$B\subseteq N'$，$\mu^*(N')=0$，则 $\mu^*(B)=0$，$B\in\cM$。所以 $A\cup B\in\cM$。

\emph{$\cM\subseteq\overline{\sigma(\cA)}$。}设 $E\in\cM$，由 C5) 写 $E=G\setminus N$，$\mu^*(N)=0$。对 $N$ 用 C3)（$\mu^*(N)=0<\infty$），得 $G'\in\sigma(\cA)$，$G'\supseteq N$，$\mu^*(G')=0$。于是
\[
E=(G\setminus G')\cup\bigl((G\cap G')\setminus N\bigr),
\]
第一部分属于 $\sigma(\cA)$，第二部分含于零测集 $G'$，属于 $\cN$。所以 $E\in\overline{\sigma(\cA)}$。

\emph{测度相同。}对 $E=A\cup B$（$A\in\sigma(\cA)$，$B\in\cN$），完备化测度的值是 $\mu^*(A)$；而 $\mu^*(A)\leqslant\mu^*(E)\leqslant\mu^*(A)+\mu^*(B)=\mu^*(A)$。
\end{solution}
\item 应用：设 $E\subseteq\R$ 是 Lebesgue 可测集，$\lambda(E)<\infty$。证明对任意 $\varepsilon>0$，存在有限个区间的并 $B$，使 $\lambda(E\triangle B)<\varepsilon$。
\begin{solution}
取 $\cA=\overline{\cS_1}$，长度是它上面的 $\sigma$-有限测度（$\R=\bigcup_n(-n,n]$），Lebesgue 可测集就是 $\cM(\lambda^*)$。由 C2)（$\varepsilon$ 换成 $\varepsilon/3$），存在 $B\in\overline{\cS_1}$ 使 $\lambda^*(E\triangle B)\leqslant\frac23\varepsilon<\varepsilon$。$B$ 是有限个左开右闭区间的不交并；由于 $\lambda(B)\leqslant\lambda(E)+\lambda(E\triangle B)<\infty$，这些区间都有界。

这个结论在分析中常用：一个有限测度的可测集，在“对称差的测度”意义下可以用有限个区间任意逼近。比如证明 $\lambda(E\cap(E+x))\to\lambda(E)$（$x\to0$）时，先对区间验证，再用它逼近。
\end{solution}
\end{questions}
\end{probleme}

\begin{probleme}{D}{不可测集}
本题证明：\emph{不存在定义在 $[0,1)$ 的全体子集上、可数可加、平移不变（按模 $1$ 平移）、且 $[0,1)$ 的测度为 $1$ 的函数。}特别地，Lebesgue 测度不能延拓到 $2^{[0,1)}$ 上并保持平移不变性。对 $x\in[0,1)$ 与 $E\subseteq[0,1)$，令 $x\oplus E=\{(x+y)\bmod1: y\in E\}$。
\begin{questions}
\item 设 $E\subseteq[0,1)$ 是 Borel 集，$x\in(0,1)$。令 $A=E\cap[0,1-x)$，$B=E\cap[1-x,1)$。证明 $x\oplus E=(A+x)\cup(B+x-1)$ 是不交并，并由平移不变性证明 $\lambda(x\oplus E)=\lambda(E)$。
\begin{solution}
若 $y\in A$，则 $y+x\in[x,1)$，$(x+y)\bmod1=y+x$；若 $y\in B$，则 $y+x\in[1,1+x)$，$(x+y)\bmod1=y+x-1\in[0,x)$。所以 $x\oplus E=(A+x)\cup(B+x-1)$，前者含于 $[x,1)$，后者含于 $[0,x)$，不交。$A,B$ 是 Borel 集，平移后仍是 Borel 集，由平移不变性（命题~\ref{prop:extension:lebesgue}(2)），
\[
\lambda(x\oplus E)=\lambda(A+x)+\lambda(B+x-1)=\lambda(A)+\lambda(B)=\lambda(E).
\]
\end{solution}
\item 在 $[0,1)$ 上定义 $x\sim y$ 当且仅当 $x-y\in\Q$。证明这是等价关系。由选择公理，取集合 $V$，使它恰好含有每个等价类中的一个元素。
\begin{solution}
自反：$x-x=0\in\Q$。对称：$x-y\in\Q$ 则 $y-x\in\Q$。传递：$x-y,\ y-z\in\Q$ 则 $x-z=(x-y)+(y-z)\in\Q$。等价类两两不交、并为 $[0,1)$。选择公理允许我们从每个等价类中取出一个元素，组成 $V$。
\end{solution}
\item 证明 $\{q\oplus V: q\in\Q\cap[0,1)\}$ 两两不交，并且它们的并是 $[0,1)$。
\begin{solution}
\emph{不交。}设 $z\in(q\oplus V)\cap(q'\oplus V)$，$z\equiv v+q\equiv v'+q'\pmod1$，$v,v'\in V$。则 $v-v'\equiv q'-q\pmod1$，即 $v-v'\in\Q$，$v\sim v'$。$V$ 中每个等价类只有一个元素，所以 $v=v'$，于是 $q\equiv q'\pmod1$，又 $q,q'\in[0,1)$，得 $q=q'$。

\emph{覆盖。}任取 $x\in[0,1)$，设 $v\in V$ 是 $x$ 所在等价类的代表元，$r=x-v\in\Q\cap(-1,1)$。若 $r\geqslant0$，取 $q=r$；若 $r<0$，取 $q=r+1\in(0,1)$。两种情形都有 $(v+q)\bmod1=x$，即 $x\in q\oplus V$。
\end{solution}
\item 假设 $\nu$ 是 $2^{[0,1)}$ 上可数可加、满足 $\nu(x\oplus E)=\nu(E)$、$\nu([0,1))=1$ 的函数。由 D3) 推出矛盾。
\begin{solution}
$\Q\cap[0,1)$ 可数，由 D3) 与可数可加性，
\[
1=\nu([0,1))=\sum_{q\in\Q\cap[0,1)}\nu(q\oplus V)=\sum_{q}\nu(V).
\]
右边是可数个相同的非负数之和：$\nu(V)=0$ 时为 $0$，$\nu(V)>0$ 时为 $\infty$，都不等于 $1$。矛盾。
\end{solution}
\item 由 D1)、D4) 的论证说明 $V$ 不是 Lebesgue 可测的。（需要把 D1) 推广到 Lebesgue 可测集：利用作业 C。）
\begin{solution}
\emph{第一步：D1) 对 Lebesgue 可测集成立。}设 $E\subseteq[0,1)$ Lebesgue 可测。由 C5)，$E=G\setminus N$，$G$ 是 Borel 集（可以换成 $G\cap[0,1)$），$\lambda^*(N)=0$；由 C3)，存在 Borel 集 $G'\supseteq N$，$\lambda(G')=0$（同样可以与 $[0,1)$ 取交）。$y\mapsto(x+y)\bmod1$ 是 $[0,1)$ 到自身的双射，所以 $x\oplus E=(x\oplus G)\setminus(x\oplus N)$，而 $x\oplus N\subseteq x\oplus G'$，由 D1) 后者是测度为 $0$ 的 Borel 集。于是 $x\oplus N$ 是 Lebesgue 可测的零测集，$x\oplus E$ Lebesgue 可测，且 $\lambda(x\oplus E)=\lambda(x\oplus G)=\lambda(G)=\lambda(E)$。

\emph{第二步。}若 $V$ Lebesgue 可测，由第一步，每个 $q\oplus V$ 可测且测度等于 $\lambda(V)$，再由 D3) 与可数可加性，$1=\lambda([0,1))=\sum_q\lambda(V)$，与 D4) 同样矛盾。所以 $V$ 不是 Lebesgue 可测的，当然也不是 Borel 集。
\end{solution}
\end{questions}
\end{probleme}

\thinkingpart

\begin{probleme}*[问题]{E}{Lebesgue 可测集比 Borel 集多}
\begin{questions}
\item 利用 Cantor 集 $C$（习题~\exref{ex:extension:A4}）和 Carathéodory 定理中“零外测度集可测”的结论，证明 $C$ 的每个子集都是 Lebesgue 可测的，从而 Lebesgue 可测集的全体的基数为 $\abs{2^\R}$。
\begin{solution}
$\lambda^*(C)=\lambda(C)=0$（习题 A4-1)），若 $D\subseteq C$，由单调性 $\lambda^*(D)=0$，由定理~\ref{thm:extension:caratheodory} 的最后一句，$D\in\cM(\lambda^*)$。所以 $\cM(\lambda^*)\supseteq2^C$，$\abs{\cM(\lambda^*)}\geqslant\abs{2^C}=\abs{2^\R}$（由 A4-3)，$\abs C=\abs\R$）；另一方面 $\cM(\lambda^*)\subseteq2^\R$。所以等号成立。
\end{solution}
\item 结合第~\ref{lec:set-systems} 讲作业 F（$\abs{\cR}=\abs\R$），证明存在不是 Borel 集的 Lebesgue 可测集。
\begin{solution}
$\abs\cR=\abs\R<\abs{2^\R}=\abs{\cM(\lambda^*)}$（Cantor 定理），所以 $\cM(\lambda^*)\not\subseteq\cR$。事实上，由 E1) 的论证，Cantor 集的子集中“绝大多数”都不是 Borel 集。
\end{solution}
\item 这个证明是非构造性的。你能否想象一个具体的例子？（第~\ref{lec:random-variables} 讲的思考题会给出一个。）
\begin{solution}
想法是用一个连续的严格递增函数 $g$ 把 Cantor 集 $C$ 映成一个测度为正的集合 $g(C)$。在 $g(C)$ 中取一个不可测集 $V'$（用作业 D 的方法），再拉回来：$M=g^{-1}(V')\subseteq C$。$M$ 是零测集的子集，所以 Lebesgue 可测；如果 $M$ 是 Borel 集，那么 $V'=g(M)$ 作为 Borel 集在连续映射 $g^{-1}$ 下的原像也是 Borel 集，矛盾。具体的 $g$ 取 $x+c(x)$，$c$ 是 Cantor 函数，细节见第~\ref{lec:random-variables} 讲问题 E。这个例子仍然用了选择公理（构造 $V'$ 时），在不用选择公理的情况下，人们写不出非 Borel 的 Lebesgue 可测集的“公式”。
\end{solution}
\end{questions}
\end{probleme}

\exercisepart{课外知识：Steinhaus 定理}

\begin{probleme}{F}{正测度集的差集含有区间}
本题证明：\emph{若 $A\in\cR$，$\lambda(A)>0$，则 $A-A=\{x-y:x,y\in A\}$ 含有一个以 $0$ 为中心的开区间。}
\begin{questions}
\item 说明不妨设 $0<\lambda(A)<\infty$。（提示：$A\cap[-n,n]$。）
\begin{solution}
$A\cap[-n,n]\uparrow A$，由从下连续，$\lambda(A\cap[-n,n])\to\lambda(A)>0$，所以某个 $A_n=A\cap[-n,n]$ 满足 $0<\lambda(A_n)\leqslant2n<\infty$。$A_n-A_n\subseteq A-A$，对 $A_n$ 证明结论即可。
\end{solution}
\item 利用习题~\exref{ex:extension:A5} 的第二部分，证明存在开集 $U\supseteq A$ 使 $\lambda(U)<\frac43\lambda(A)$。
\begin{solution}
由 A5-2)，$\lambda(A)=\lambda^*(A)=\inf\{\lambda(U):U\supseteq A\ \text{开}\}$。由于 $0<\lambda(A)<\infty$，$\frac43\lambda(A)>\lambda(A)$，由下确界的定义存在开集 $U\supseteq A$ 使 $\lambda(U)<\frac43\lambda(A)$。
\end{solution}
\item 由引理~\ref{lem:set-systems:open-sets-R}，$U$ 是可数个不交开区间 $I_k$ 的并（先证明开集是可数个\emph{不交}开区间的并：考虑包含一点的最大开区间）。证明存在某个 $k$ 使 $\lambda(A\cap I_k)>\frac34\lambda(I_k)$。
\begin{solution}
\emph{开集的结构。}对 $x\in U$，令 $I_x$ 为 $U$ 中包含 $x$ 的一切开区间之并。若干个含公共点 $x$ 的区间之并仍是区间，所以 $I_x$ 是含于 $U$ 的开区间，并且是含 $x$ 的最大者。若 $I_x\cap I_y\neq\varnothing$，则 $I_x\cup I_y$ 是含 $x$ 的开区间，由最大性 $I_y\subseteq I_x$，同理 $I_x\subseteq I_y$，所以不同的 $I_x$ 两两不交。每个 $I_x$ 含有理数，不同的 $I_x$ 含不同的有理数，所以只有可数个，记为 $I_1,I_2,\dots$。它们都有界，因为 $\lambda(U)<\infty$。

\emph{选区间。}若对一切 $k$ 都有 $\lambda(A\cap I_k)\leqslant\frac34\lambda(I_k)$，则由可数可加性
\[
\lambda(A)=\sum_k\lambda(A\cap I_k)\leqslant\frac34\sum_k\lambda(I_k)=\frac34\lambda(U)<\frac34\cdot\frac43\lambda(A)=\lambda(A),
\]
矛盾。
\end{solution}
\item 记 $I=I_k$，$\ell=\lambda(I)$，$B=A\cap I$。设 $\abs x<\ell/2$。证明 $B\cup(B+x)\subseteq I\cup(I+x)$，后者是长度小于 $\frac32\ell$ 的区间。
\begin{solution}
$B\subseteq I$，$B+x\subseteq I+x$。设 $I=(c,c+\ell)$，不妨 $x\geqslant0$，则 $I\cup(I+x)=(c,c+\ell+x)$（因为 $x<\ell$，两个区间相交），长度 $\ell+x<\frac32\ell$；$x<0$ 时同理为 $(c+x,c+\ell)$。
\end{solution}
\item 若 $B\cap(B+x)=\varnothing$，由平移不变性推出 $\lambda(B\cup(B+x))=2\lambda(B)>\frac32\ell$，与 F4) 矛盾。所以 $B\cap(B+x)\neq\varnothing$，即 $x\in B-B\subseteq A-A$。
\begin{solution}
若不交，$\lambda(B\cup(B+x))=\lambda(B)+\lambda(B+x)=2\lambda(B)>2\cdot\frac34\ell=\frac32\ell$，而由 F4) 它不超过 $\lambda(I\cup(I+x))<\frac32\ell$，矛盾。所以存在 $z\in B\cap(B+x)$，即 $z\in B$ 且 $z=b+x$，$b\in B$，于是 $x=z-b\in B-B$。这对一切 $\abs x<\ell/2$ 成立，所以 $(-\frac\ell2,\frac\ell2)\subseteq A-A$。
\end{solution}
\item 推论：若 $f:\R\to\R$ 满足 $f(x+y)=f(x)+f(y)$，并且 $f$ 是 Borel 可测的（即 $\{f\leqslant c\}\in\cR$ 对一切 $c$ 成立），则 $f(x)=f(1)x$。（提示：某个 $\{\abs f\leqslant n\}$ 有正测度，由 Steinhaus 定理 $f$ 在 $0$ 附近有界。）
\begin{solution}
\emph{第一步：$f$ 在有理数上是线性的。}由可加性，$f(0)=2f(0)$，$f(0)=0$；$f(-x)=-f(x)$；$f(nx)=nf(x)$（归纳）；$f(x)=nf(x/n)$，所以 $f(\frac mnx)=\frac mnf(x)$。特别地 $f(q)=f(1)q$，$q\in\Q$。

\emph{第二步：$f$ 在 $0$ 附近有界。}$E_n=\{\abs f\leqslant n\}=\{f\leqslant n\}\setminus\bigcup_k\{f\leqslant-n-\frac1k\}$ 是 Borel 集，$E_n\uparrow\R$，由从下连续某个 $\lambda(E_n)>0$。由 F5)，存在 $\delta>0$ 使 $(-\delta,\delta)\subseteq E_n-E_n$。对 $\abs z<\delta$，$z=x-y$，$x,y\in E_n$，$\abs{f(z)}=\abs{f(x)-f(y)}\leqslant2n$。

\emph{第三步：连续。}对 $\abs z<\delta/m$，$\abs{f(z)}=\frac1m\abs{f(mz)}\leqslant\frac{2n}m$。所以 $f$ 在 $0$ 处连续，再由 $f(x+h)-f(x)=f(h)$，$f$ 处处连续。

\emph{第四步。}连续函数 $f(x)$ 与 $f(1)x$ 在稠密集 $\Q$ 上相等，所以处处相等。

不可测的可加函数是存在的（借助选择公理，取 $\R$ 作为 $\Q$ 上线性空间的一组基），这个推论说明它们都极不规则。
\end{solution}
\end{questions}
\end{probleme}
